%% 4-7 三角形を解く：正弦定理と余弦定理（第13回：組 V4-4。ページ制限なし。偶数ページ）
%% 4-0 が「逆運動学の解き方は 4-7 で」と指している → 道具⑥（2リンクのロボットアーム。余弦定理でひじの角、順運動学の式も）
%% 物理シリーズ（力学）のために：正弦定理（3つの力のつり合い）と余弦定理（合力の大きさ）→ 道具⑦
\makeatletter
\newcommand\bkivhline[1]{\par\addvspace{3pt}{\centering #1\par}\addvspace{3pt}\noindent\ignorespaces}
\newcommand\bkivhstep[2]{\par\addvspace{6pt}\Needspace{4\baselineskip}\noindent
  {\bkcircled[bkOrange]{#1}\enspace\sffamily\bfseries\color{bkOrange!75!black}#2}\par\nobreak\vspace{2pt}}
\newcommand\bkivhhead[1]{\par\addvspace{4pt}\Needspace{5\baselineskip}\noindent\textbf{#1}\par\nobreak}
% 直角のしるし  \bkivhright{角の点}{1本めの向き(度)}{大きさ mm}
\newcommand\bkivhright[3]{\draw[bkInk, line width=0.5pt] ($#1+(#2:#3)$) -- ++(#2+90:#3) -- ++(#2+180:#3);}
% ロボットのうで（2リンク）  \bkivharm{L1(mm)}{L2(mm)}{θ1}{θ2}{色}
\newcommand\bkivharm[5]{%
  \coordinate (O) at (0,0);
  \coordinate (E) at ({#1*cos(#3)},{#1*sin(#3)});
  \coordinate (P) at ($(E)+({#2*cos(#3+#4)},{#2*sin(#3+#4)})$);
  \draw[#5, line width=3.2pt, line cap=round] (O) -- (E) -- (P);
  \fill[bkGray!70] (-4,-3) rectangle (4,0);
  \filldraw[fill=white, draw=bkInk, line width=0.7pt] (O) circle (1.6);
  \filldraw[fill=white, draw=bkInk, line width=0.7pt] (E) circle (1.4);
  \fill[bkInk] (P) circle (1.1);}
\makeatother

\bkchapter{4}{7}{三角形を解く：\\正弦定理と余弦定理}{\bklevel{3}}{三角形を解く：正弦定理と余弦定理}
\bkmichishirube{4-1,4-4,4-6}{4-9,4-10,4-14}

\begin{goal}
\begin{itemize}[label={\color{\bkTT}●}, leftmargin=1.6\zw, itemsep=0pt, topsep=1pt]
\item \bklevel{3}\enspace 鈍角の三角比を使い、三角形の面積を $\dfrac{1}{2}ab\sin C$ で求められる。
\item \bklevel{3}\enspace 正弦定理と余弦定理を使い分けて、三角形の辺や角を求められる。
\item \bklevel{3}\enspace 2つの関節をもつロボットのうでの角を、余弦定理で求められる。
\end{itemize}
\end{goal}

\begin{hatena}
\noindent\begin{minipage}[c]{0.56\linewidth}
川の向こう岸に、大きな木Cがあります。木までの距離を知りたいのですが、川は渡れません。\par\vspace{2pt}
そこで、こちらの岸に2つの点A、Bをとり、ABの長さと、AとBから木を見た角をはかりました。$\mathrm{AB}=\qty{100}{m}$、Aの角は\ang{60}、Bの角は\ang{75}です。\par\vspace{2pt}
直角のない三角形でも、AからCまでの距離は計算できるでしょうか。
\end{minipage}\hfill
\begin{minipage}[c]{0.4\linewidth}
\centering
\begin{tikzpicture}[x=0.9mm, y=0.9mm]
  \fill[bkBlue!20] (-6,7) rectangle (44,29);
  \draw[bkBlue!60, line width=0.6pt] (-6,7) -- (44,7) (-6,29) -- (44,29);
  \coordinate (A) at (0,0); \coordinate (B) at (30,0);
  \coordinate (C) at ({40.98*cos(60)},{40.98*sin(60)});
  \fill[bkStem] ($(C)+(-0.9,0)$) rectangle ($(C)+(0.9,3)$);
  \fill[bkLeaf] ($(C)+(0,5.5)$) circle (3.2);
  \draw[bkInk, line width=0.9pt] (A) -- (B);
  \draw[bkOrange, line width=0.9pt] (B) -- (C);
  \draw[bkRed, line width=1.1pt, dashed] (A) -- (C);
  \fill[bkInk] (A) circle (0.8) (B) circle (0.8) (C) circle (0.8);
  \node[font=\sffamily\footnotesize, anchor=north] at ($(A)+(0,-1)$) {A};
  \node[font=\sffamily\footnotesize, anchor=north] at ($(B)+(0,-1)$) {B};
  \node[font=\sffamily\footnotesize, anchor=east] at ($(C)+(-3.6,4)$) {C};
  \node[font=\sffamily\footnotesize, anchor=north] at (15,-1.2) {\qty{100}{m}};
  \draw[bkInk, line width=0.5pt] (5,0) arc (0:60:5);
  \node[font=\sffamily\footnotesize] at (30:8) {\ang{60}};
  \draw[bkInk, line width=0.5pt] ($(B)+(-5,0)$) arc (180:105:5);
  \node[font=\sffamily\footnotesize] at ($(B)+(140:8.2)$) {\ang{75}};
  \node[font=\sffamily\footnotesize\bfseries, text=bkRed!80!black, fill=bkBlue!20, inner sep=0.8pt] at ($(A)!0.55!(C)+(-3.5,0)$) {？};
  \node[font=\sffamily\footnotesize, text=bkBlue!70!black, anchor=east] at (44,18) {川};
\end{tikzpicture}
\end{minipage}
\end{hatena}

\noindent 4-6では、直角三角形の辺の比を使いました。この章では、直角のない三角形でも、辺と角の関係を式にします。三角形のわかっている辺や角から、のこりの辺や角を求めることを、\textbf{三角形を解く}といいます。

\begin{yattemiyou}[辺と、向かいの角の $\sin$ をくらべよう]
\noindent\begin{minipage}[c]{0.6\linewidth}
\textbf{用意するもの}\par
ノート、定規、分度器、電卓
\begin{enumerate}[label={\bkcircled[bkGreen]{\arabic*}}, leftmargin=2.4\zw, itemsep=1pt, topsep=3pt]
\item $\mathrm{AB}=\qty{8}{cm}$、Aの角\ang{50}、Bの角\ang{70}の三角形ABCをかこう。
\item 辺BC、CAの長さと、Cの角をはかろう。
\item 「辺の長さ÷向かいの角の $\sin$」を、3つの辺について計算しよう。
\item $\mathrm{AB}=\qty{6}{cm}$、Aの角\ang{30}、Bの角\ang{110}の三角形でもやってみよう。
\end{enumerate}
\end{minipage}\hfill
\begin{minipage}[c]{0.37\linewidth}\centering
\begin{tikzpicture}[x=0.88mm, y=0.88mm]
  \coordinate (A) at (0,0); \coordinate (B) at (40,0);
  \coordinate (C) at ({43.4*cos(50)},{43.4*sin(50)});
  \fill[bkGreen!10] (A) -- (B) -- (C) -- cycle;
  \draw[bkGreen!60!black, line width=0.9pt] (A) -- (B) -- (C) -- cycle;
  \node[font=\sffamily\footnotesize, anchor=north east] at (A) {A};
  \node[font=\sffamily\footnotesize, anchor=north west] at (B) {B};
  \node[font=\sffamily\footnotesize, anchor=south] at ($(C)+(0,0.6)$) {C};
  \draw[bkInk, line width=0.5pt] (6,0) arc (0:50:6);
  \node[font=\sffamily\footnotesize] at (25:9.5) {\ang{50}};
  \draw[bkInk, line width=0.5pt] ($(B)+(-6,0)$) arc (180:110:6);
  \node[font=\sffamily\footnotesize] at ($(B)+(145:9.5)$) {\ang{70}};
  \node[font=\sffamily\footnotesize, anchor=north] at (20,-1.2) {\qty{8}{cm}};
\end{tikzpicture}
\end{minipage}
\par\vspace{4pt}
ミナさんがはかった結果です（長さは\qty{0.1}{cm}まで）。
\begin{bkcenter}
\begin{bktable}{colspec={Q[c,m,wd=28mm] Q[c,m,wd=44mm] Q[c,m,wd=44mm]}, colsep=3pt, rowsep=1.2pt, row{1}={font=\sffamily\footnotesize}}
 & ①の三角形 & ④の三角形 \\
BC÷$\sin A$ & $7.1\div\sin\ang{50}=9.27$ & $4.7\div\sin\ang{30}=9.40$ \\
CA÷$\sin B$ & $8.7\div\sin\ang{70}=9.26$ & $8.8\div\sin\ang{110}=9.36$ \\
AB÷$\sin C$ & $8\div\sin\ang{60}=9.24$ & $6\div\sin\ang{40}=9.33$ \\
\end{bktable}
\end{bkcenter}
\end{yattemiyou}

\noindent どちらの三角形でも、3つの値がほとんど同じになりました（小さなずれは、はかるときのずれ）。Cの角は、内角の和が\ang{180}なので、①は $\ang{180}-\ang{50}-\ang{70}=\ang{60}$、④は\ang{40}です。\ang{110}のような鈍角の $\sin$ も、4-6⑧の決め方で計算しています。この「同じになる値」のなぞを、③でときます。

\begin{omoidasou}
\bkrecall{4-1}{三角形の内角の和は\ang{180}。合同条件は「3辺」「2辺とそのあいだの角」「1辺とその両はしの角」で、どれも三角形が1つに決まる条件。}
\bkrecall{4-2}{円周角は中心角の半分。同じ弧に対する円周角は等しい。直径に対する円周角は\ang{90}。三角形の3つの頂点を通る円の中心が外心。}
\bkrecall{4-4}{2点 $(x_1,\ y_1)$、$(x_2,\ y_2)$ の距離は $\sqrt{(x_2-x_1)^2+(y_2-y_1)^2}$。}
\bkrecall{4-6}{対辺＝斜辺×$\sin\theta$、隣辺＝斜辺×$\cos\theta$。$\sin^2\theta+\cos^2\theta=1$。\\半径 $r$ の半円の上で角 $\theta$ の点は $(r\cos\theta,\ r\sin\theta)$。\\$\sin(\ang{180}-\theta)=\sin\theta$、$\cos(\ang{180}-\theta)=-\cos\theta$（鈍角の $\cos$ は負）。}
\end{omoidasou}

%% ---------------- 数学の道具 ----------------
\bksection{doguu}{数学の道具}

\bksubsection{1}{三角形の名前のつけ方と、鈍角の三角比}

\leavevmode\label{bk:4-7:donkaku}この章では、三角形ABCの辺と角を、次のように書きます。

\noindent\begin{minipage}[c]{0.56\linewidth}
\begin{itemize}[leftmargin=1.2\zw, itemsep=2pt, topsep=0pt, label={\color{\bkTT}●}]
\item 頂点Aの角（$\angle\mathrm{BAC}$。4-1）を、かんたんに $A$ と書く。$B$、$C$ も同じ。
\item 角 $A$ の向かいの辺BCの長さを $a$、角 $B$ の向かいの辺CAを $b$、角 $C$ の向かいの辺ABを $c$ と書く。
\end{itemize}
\par\vspace{2pt}
大文字の角と、同じ文字の小文字の辺は、いつも向かい合っています。4-6で「$\theta$ の対辺」とよんだ辺と同じ見方です。
\end{minipage}\hfill
\begin{minipage}[c]{0.4\linewidth}\centering
\begin{tikzpicture}[x=1mm, y=1mm]
  \coordinate (A) at (0,0); \coordinate (B) at (46,0); \coordinate (C) at (16,26);
  \fill[bkTheme!12] (A) -- (B) -- (C) -- cycle;
  \draw[bkInk!70, line width=0.8pt] (A) -- (B);
  \draw[bkRed, line width=1.5pt] (B) -- (C);
  \draw[bkBlue, line width=1.5pt] (C) -- (A);
  \draw[bkGreen!70!black, line width=1.5pt] (A) -- (B);
  \node[font=\sffamily\footnotesize, anchor=north east] at (A) {A};
  \node[font=\sffamily\footnotesize, anchor=north west] at (B) {B};
  \node[font=\sffamily\footnotesize, anchor=south] at ($(C)+(0,0.5)$) {C};
  \node[font=\sffamily\footnotesize, text=bkRed!80!black, anchor=south west] at ($(B)!0.5!(C)+(0.5,0.5)$) {$a$};
  \node[font=\sffamily\footnotesize, text=bkBlue!80!black, anchor=south east] at ($(A)!0.5!(C)+(-0.5,0.5)$) {$b$};
  \node[font=\sffamily\footnotesize, text=bkGreen!50!black, anchor=north] at (23,-1.2) {$c$};
  \draw[bkInk, line width=0.5pt] (5,0) arc (0:58.4:5);
  \node[font=\sffamily\footnotesize] at (29:7.8) {$A$};
\end{tikzpicture}
\end{minipage}

\bkivhstep{1}{三角形を座標の上に置く}
三角形の頂点Cを原点に、辺CBを $x$ 軸の正の向きに置きます。Bの座標は $(a,\ 0)$ です。Aは、Cから長さ $b$ で、$x$ 軸から角 $C$ だけ回った所にあります。半径 $b$ の半円の上の点なので（4-6⑧）、
\bkivhline{\fbox{$\mathrm{A}(b\cos C,\ b\sin C)$}}
です。この置き方は、$C$ が鋭角でも鈍角でも使えます。

\begin{bkfigure}
\begin{tikzpicture}[x=1mm, y=1mm]
  \begin{scope}[shift={(8,0)}]
  \draw[bkGray, line width=0.6pt, -{Stealth[length=4pt]}] (-6,0) -- (50,0) node[anchor=north, font=\sffamily\footnotesize] {$x$};
  \coordinate (C) at (0,0); \coordinate (B) at (40,0); \coordinate (A) at ({30*cos(60)},{30*sin(60)});
  \fill[bkTheme!12] (A) -- (B) -- (C) -- cycle;
  \draw[bkInk, line width=0.9pt] (A) -- (B) -- (C) -- cycle;
  \draw[bkOrange, line width=1.3pt] (A) -- ({30*cos(60)},0);
  \draw[bkBlue, line width=1.3pt] (C) -- ({30*cos(60)},0);
  \bkivhright{({30*cos(60)},0)}{90}{2}
  \draw[bkInk, line width=0.5pt] (6,0) arc (0:60:6);
  \node[font=\sffamily\footnotesize] at (30:9) {$C$};
  \node[font=\sffamily\footnotesize, anchor=north east] at (-0.6,-0.6) {C};
  \node[font=\sffamily\footnotesize, anchor=north] at (40,-1.2) {B};
  \node[font=\sffamily\footnotesize, anchor=south] at ($(A)+(0,0.6)$) {A};
  \node[font=\sffamily\footnotesize, anchor=south east] at ($(C)!0.5!(A)+(-0.5,0)$) {$b$};
  \node[font=\sffamily\footnotesize, text=bkBlue!80!black, anchor=north] at (7.5,-1.2) {$b\cos C$};
  \node[font=\sffamily\footnotesize, text=bkOrange!80!black, anchor=west] at ({30*cos(60)+1},11) {$b\sin C$};
  \node[font=\sffamily\small, anchor=north] at (20,-7.5) {$C$ が鋭角（$\cos C>0$）};
  \end{scope}
  \begin{scope}[shift={(96,0)}]
  \draw[bkGray, line width=0.6pt, -{Stealth[length=4pt]}] (-26,0) -- (34,0) node[anchor=north, font=\sffamily\footnotesize] {$x$};
  \coordinate (C) at (0,0); \coordinate (B) at (26,0); \coordinate (A) at ({30*cos(120)},{30*sin(120)});
  \fill[bkTheme!12] (A) -- (B) -- (C) -- cycle;
  \draw[bkInk, line width=0.9pt] (A) -- (B) -- (C) -- cycle;
  \draw[bkOrange, line width=1.3pt] (A) -- ({30*cos(120)},0);
  \draw[bkBlue, line width=1.3pt] (C) -- ({30*cos(120)},0);
  \bkivhright{({30*cos(120)},0)}{0}{2}
  \draw[bkInk, line width=0.5pt] (5,0) arc (0:120:5);
  \node[font=\sffamily\footnotesize] at (60:7.5) {$C$};
  \node[font=\sffamily\footnotesize, anchor=north west] at (0.6,-0.6) {C};
  \node[font=\sffamily\footnotesize, anchor=north] at (26,-1.2) {B};
  \node[font=\sffamily\footnotesize, anchor=south] at ($(A)+(0,0.6)$) {A};
  \node[font=\sffamily\footnotesize, text=bkBlue!80!black, anchor=north east] at (-1.5,-1.2) {$b\cos C<0$};
  \node[font=\sffamily\footnotesize, text=bkOrange!80!black, anchor=east] at ({30*cos(120)-1},11) {$b\sin C$};
  \node[font=\sffamily\small, anchor=north] at (4,-7.5) {$C$ が鈍角（$\cos C<0$）};
  \end{scope}
\end{tikzpicture}
\bkfigcaption{Aから $x$ 軸へ下ろした垂線の足は、$C$ が鈍角だと、Cより左（Bの反対がわ）に出る。}
\end{bkfigure}

\bkivhstep{2}{鈍角の $\cos$ が負になるわけを、図で見る}
$C$ が鈍角のとき、Aの $x$ 座標 $b\cos C$ は負です。図では、Aから下ろした垂線の足が、Cから見てBと反対がわに出ています。「$\cos C$ が負」は、「Aの真下の点が、Cの向こうがわにある」という意味なのです。高さ $b\sin C$ は、鋭角でも鈍角でも正です。数で確かめます。$b=4$ とすると、
\begin{itemize}[leftmargin=1.2\zw, itemsep=1pt, topsep=2pt, label={\color{\bkTT}●}]
\item $C=\ang{60}$：$\mathrm{A}(4\cos\ang{60},\ 4\sin\ang{60})=(2,\ 3.46)$
\item $C=\ang{120}$：$\mathrm{A}(4\cos\ang{120},\ 4\sin\ang{120})=(-2,\ 3.46)$
\end{itemize}
\ang{60}と\ang{120}で、高さは同じで、横の位置だけが左右反対になります。これは、$\cos\ang{120}=-\cos\ang{60}=-0.5$ だからです。

\bksubsection{2}{三角形の面積：$S=\dfrac{1}{2}ab\sin C$}

\leavevmode\label{bk:4-7:menseki}\noindent\begin{minipage}[c]{0.58\linewidth}
三角形の面積は、底辺×高さ÷2でした（4-5）。①の置き方で、底辺をCB（長さ $a$）にすると、高さはAの $y$ 座標 $b\sin C$ です。だから、
\bkivhline{\fbox{$S=\dfrac{1}{2}\times a\times b\sin C=\dfrac{1}{2}ab\sin C$}}
\textbf{2辺と、そのあいだの角}がわかれば、面積が求められます。
\end{minipage}\hfill
\begin{minipage}[c]{0.38\linewidth}\centering
\begin{tikzpicture}[x=0.9mm, y=0.9mm]
  \coordinate (C) at (0,0); \coordinate (B) at (44,0); \coordinate (A) at ({30*cos(50)},{30*sin(50)});
  \fill[bkTheme!14] (A) -- (B) -- (C) -- cycle;
  \draw[bkInk, line width=0.9pt] (A) -- (B) -- (C) -- cycle;
  \draw[bkOrange, line width=1.2pt, dashed] (A) -- ({30*cos(50)},0);
  \bkivhright{({30*cos(50)},0)}{90}{2}
  \draw[bkInk, line width=0.5pt] (6,0) arc (0:50:6);
  \node[font=\sffamily\footnotesize] at (25:8.8) {$C$};
  \node[font=\sffamily\footnotesize, anchor=north] at (22,-1.2) {$a$（底辺）};
  \node[font=\sffamily\footnotesize, anchor=south east] at ($(C)!0.55!(A)+(-0.4,0)$) {$b$};
  \node[font=\sffamily\footnotesize, text=bkOrange!80!black, anchor=west] at ({30*cos(50)+0.8},9) {$b\sin C$};
\end{tikzpicture}
\end{minipage}

\par\vspace{4pt}
\textbf{確かめ1：直角のとき}\quad $C=\ang{90}$ なら $\sin\ang{90}=1$ で、$S=\dfrac{1}{2}ab$。直角をはさむ2辺の積の半分で、長方形の半分と合います。

\textbf{確かめ2：数で}\quad $a=6$、$b=4$、$C=\ang{30}$ なら、高さは $4\sin\ang{30}=2$ なので、面積は $6\times2\div2=6$。式でも $\dfrac{1}{2}\times6\times4\times\dfrac{1}{2}=6$ です。

\noindent\begin{minipage}[c]{0.5\linewidth}\centering
\begin{tikzpicture}[x=1mm, y=1mm]
  \coordinate (C) at (0,0); \coordinate (B) at (30,0);
  \coordinate (A1) at ({24*cos(60)},{24*sin(60)});
  \coordinate (A2) at ({24*cos(120)},{24*sin(120)});
  \fill[bkBlue!14] (A1) -- (B) -- (C) -- cycle;
  \fill[bkOrange!14] (A2) -- (B) -- (C) -- cycle;
  \draw[bkBlue, line width=0.9pt] (A1) -- (B) -- (C) -- cycle;
  \draw[bkOrange, line width=0.9pt] (A2) -- (B) (C) -- (A2);
  \draw[bkGray, line width=0.5pt, dashed] (-20,{24*sin(60)}) -- (32,{24*sin(60)});
  \draw[bkBlue, line width=0.5pt] (4,0) arc (0:60:4);
  \node[font=\sffamily\footnotesize, text=bkBlue!80!black] at (24:7.4) {\ang{60}};
  \draw[bkOrange, line width=0.5pt] (8.5,0) arc (0:120:8.5);
  \node[font=\sffamily\footnotesize, text=bkOrange!80!black] at (92:12.5) {\ang{120}};
  \node[font=\sffamily\footnotesize, anchor=north] at (0,-1.2) {C};
  \node[font=\sffamily\footnotesize, anchor=north] at (30,-1.2) {B};
\end{tikzpicture}
\end{minipage}\hfill
\begin{minipage}[c]{0.46\linewidth}
2辺の長さが同じで、あいだの角が\ang{60}の三角形（青）と\ang{120}の三角形（オレンジ）は、高さが同じなので、面積も同じです。$\sin\ang{120}=\sin\ang{60}$ だからです（4-6⑧）。
\end{minipage}

\par\vspace{4pt}
どの2辺を使ってもかまいません。2辺 $b$、$c$ とそのあいだの角 $A$ なら $S=\dfrac{1}{2}bc\sin A$、2辺 $c$、$a$ とそのあいだの角 $B$ なら $S=\dfrac{1}{2}ca\sin B$ です。この三角形を2つ合わせた平行四辺形の面積は $ab\sin C$ です。この形は、4-14の外積で、もう一度出てきます。

\bksubsection{3}{正弦定理}

\leavevmode\label{bk:4-7:seigen}\bkyattemiyouref で、「辺÷向かいの角の $\sin$」が、3つとも同じになりました。そのわけを、2通りで考えます。

\bkivhstep{1}{高さを2通りで表す}
\noindent\begin{minipage}[c]{0.56\linewidth}
CからABへ垂直な線を下ろし、足をHとします。CHの長さ $h$ を、2つの直角三角形で表します。
\begin{itemize}[leftmargin=1.2\zw, itemsep=1pt, topsep=2pt, label={\color{\bkTT}●}]
\item 三角形ACH：斜辺 $b$、角 $A$ の対辺が $h$。だから $h=b\sin A$。
\item 三角形BCH：斜辺 $a$、角 $B$ の対辺が $h$。だから $h=a\sin B$。
\end{itemize}
同じ $h$ なので、$a\sin B=b\sin A$。両辺を $\sin A\sin B$ でわると（2-14）、
\end{minipage}\hfill
\begin{minipage}[c]{0.4\linewidth}\centering
\begin{tikzpicture}[x=0.9mm, y=0.9mm]
  \coordinate (A) at (0,0); \coordinate (B) at (50,0); \coordinate (C) at (18,28); \coordinate (H) at (18,0);
  \fill[bkBlue!12] (A) -- (H) -- (C) -- cycle;
  \fill[bkOrange!12] (H) -- (B) -- (C) -- cycle;
  \draw[bkInk, line width=0.9pt] (A) -- (B) -- (C) -- cycle;
  \draw[bkInk, line width=0.9pt, dashed] (C) -- (H);
  \bkivhright{(H)}{0}{2}
  \node[font=\sffamily\footnotesize, anchor=north east] at (A) {A};
  \node[font=\sffamily\footnotesize, anchor=north west] at (B) {B};
  \node[font=\sffamily\footnotesize, anchor=south] at ($(C)+(0,0.5)$) {C};
  \node[font=\sffamily\footnotesize, anchor=north] at ($(H)+(0,-1)$) {H};
  \node[font=\sffamily\footnotesize, text=bkBlue!80!black, anchor=south east] at ($(A)!0.5!(C)+(-0.4,0)$) {$b$};
  \node[font=\sffamily\footnotesize, text=bkOrange!80!black, anchor=south west] at ($(B)!0.5!(C)+(0.4,0)$) {$a$};
  \node[font=\sffamily\footnotesize, anchor=west] at (18.8,13) {$h$};
  \draw[bkInk, line width=0.5pt] (5,0) arc (0:57.3:5);
  \node[font=\sffamily\footnotesize] at (28:7.6) {$A$};
  \draw[bkInk, line width=0.5pt] ($(B)+(-5,0)$) arc (180:138.8:5);
  \node[font=\sffamily\footnotesize] at ($(B)+(160:7.8)$) {$B$};
\end{tikzpicture}
\end{minipage}
\bkivhline{$\dfrac{a\sin B}{\sin A\sin B}=\dfrac{b\sin A}{\sin A\sin B}$\qquad つまり\qquad $\dfrac{a}{\sin A}=\dfrac{b}{\sin B}$}
AからBCへ垂直な線を下ろして同じように考えると、$\dfrac{b}{\sin B}=\dfrac{c}{\sin C}$ も出てきます。だから、3つとも等しいのです。

\begin{seikaku}
\textbf{$A$ が鈍角のとき}\par
$A$ が鈍角なら、Hは辺ABをのばした先（Aの外がわ）に出ます。そのとき三角形ACHの、Aの所の角は $\ang{180}-A$ なので、$h=b\sin(\ang{180}-A)$。$\sin(\ang{180}-A)=\sin A$（4-6⑧）なので、やはり $h=b\sin A$ です。だから、鈍角があっても同じ式になります。
\end{seikaku}

\bkivhstep{2}{その値は、外接円の直径}
\noindent\begin{minipage}[c]{0.56\linewidth}
三角形の3つの頂点を通る円を\textbf{外接円}（がいせつえん）といいます（中心は外心。4-2）。半径を $R$ とします。\bkidx{がいせつえん}{外接円}\par\vspace{2pt}
Bから円の中心Oを通る直径BDを引きます。
\begin{enumerate}[label={\bkcircled[bkOrange]{\arabic*}}, leftmargin=2.4\zw, itemsep=1pt, topsep=2pt]
\item Dの角 $\angle\mathrm{BDC}$ と、Aの角 $A$ は、同じ弧BCに対する円周角なので等しい（4-2）。
\item BDは直径なので、$\angle\mathrm{BCD}=\ang{90}$（4-2）。
\item 直角三角形BCDで、斜辺BDは $2R$、Dの対辺はBC＝$a$。だから $a=2R\sin D=2R\sin A$。
\end{enumerate}
\end{minipage}\hfill
\begin{minipage}[c]{0.4\linewidth}\centering
\begin{tikzpicture}[x=1mm, y=1mm]
  \def\R{23}
  \coordinate (O) at (0,0);
  \coordinate (B) at (210:\R); \coordinate (C) at (330:\R); \coordinate (A) at (105:\R); \coordinate (D) at (30:\R);
  \draw[bkGray, line width=0.8pt] (O) circle (\R);
  \fill[bkTheme!12] (A) -- (B) -- (C) -- cycle;
  \draw[bkInk, line width=0.9pt] (A) -- (B) -- (C) -- cycle;
  \draw[bkBlue, line width=1.1pt] (B) -- (D);
  \draw[bkBlue, line width=1.1pt] (C) -- (D);
  \fill[bkInk] (O) circle (0.7);
  \node[font=\sffamily\footnotesize, anchor=south] at ($(O)+(0,0.6)$) {O};
  \bkivhright{(C)}{90}{2}
  \node[font=\sffamily\footnotesize, anchor=north east] at (B) {B};
  \node[font=\sffamily\footnotesize, anchor=north west] at (C) {C};
  \node[font=\sffamily\footnotesize, anchor=south] at ($(A)+(0,0.5)$) {A};
  \node[font=\sffamily\footnotesize, anchor=south west] at (D) {D};
  \draw[bkInk, line width=0.5pt] ($(A)+(-112.5:5)$) arc (-112.5:-52.5:5);
  \node[font=\sffamily\footnotesize] at ($(A)+(-82.5:7.6)$) {$A$};
  \draw[bkBlue, line width=0.5pt] ($(D)+(210:5)$) arc (210:270:5);
  \node[font=\sffamily\footnotesize, text=bkBlue!80!black] at ($(D)+(240:7.8)$) {$A$};
  \node[font=\sffamily\footnotesize, text=bkBlue!80!black, anchor=north west] at ($(O)!0.5!(B)+(0.6,-1)$) {$2R$};
  \node[font=\sffamily\footnotesize, anchor=north] at ($(B)!0.5!(C)+(0,-0.8)$) {$a$};
\end{tikzpicture}
\end{minipage}
\par\vspace{4pt}
③の式の両辺を $\sin A$ でわると $\dfrac{a}{\sin A}=2R$ です。①で見たように、ほかの辺も同じ値なので、次の式が成り立ちます。これを\textbf{正弦定理}（せいげんていり）といいます（定理。4-1）。\bkidx{せいげんていり}{正弦定理}
\bkivhline{\fbox{$\dfrac{a}{\sin A}=\dfrac{b}{\sin B}=\dfrac{c}{\sin C}=2R$}}
\bkyattemiyouref の①の三角形では、値が約9.24でした。外接円の直径が約\qty{9.24}{cm}、半径が約\qty{4.6}{cm}ということです。コンパスで、半径\qty{4.6}{cm}の円をかいて確かめてみましょう。3つの頂点の近くを通るはずです。

\bkivhhead{正弦定理が言っていること}
$a:b:c=\sin A:\sin B:\sin C$ です。\textbf{大きい角の向かいには、長い辺}があります。ただし、辺の長さは角の大きさに比例するのではなく、角の $\sin$ に比例します。たとえば、三角定規（\ang{30}・\ang{60}・\ang{90}）の辺の比は、$1:2:3$ ではなく、$\sin\ang{30}:\sin\ang{60}:\sin\ang{90}=\dfrac{1}{2}:\dfrac{\sqrt3}{2}:1=1:\sqrt3:2$ です（4-4と同じ）。

\bkivhhead{\bkhatenaref の答え}
\noindent\begin{minipage}[c]{0.6\linewidth}
三角形ABCで、$c=\mathrm{AB}=100$、$A=\ang{60}$、$B=\ang{75}$ です。まず、内角の和から
\bkivhline{$C=\ang{180}-\ang{60}-\ang{75}=\ang{45}$}
求めたいACは、角 $B$ の向かいの辺 $b$ です。正弦定理 $\dfrac{b}{\sin B}=\dfrac{c}{\sin C}$ の両辺に $\sin B$ をかけて、
\end{minipage}\hfill
\begin{minipage}[c]{0.36\linewidth}\centering
\begin{tikzpicture}[x=1mm, y=1mm]
  \coordinate (A) at (0,0); \coordinate (B) at (30,0);
  \coordinate (C) at ({40.98*cos(60)},{40.98*sin(60)});
  \fill[bkTheme!12] (A) -- (B) -- (C) -- cycle;
  \draw[bkInk, line width=0.8pt] (A) -- (B) -- (C) -- cycle;
  \draw[bkRed, line width=1.3pt] (A) -- (C);
  \node[font=\sffamily\footnotesize, anchor=north] at (A) {A};
  \node[font=\sffamily\footnotesize, anchor=north] at (B) {B};
  \node[font=\sffamily\footnotesize, anchor=south] at ($(C)+(0,0.5)$) {C};
  \node[font=\sffamily\footnotesize, anchor=north] at (15,-1) {100};
  \node[font=\sffamily\footnotesize, text=bkRed!80!black, anchor=east] at ($(A)!0.5!(C)+(-0.6,0)$) {$b$};
  \draw[bkInk, line width=0.5pt] (4,0) arc (0:60:4);
  \node[font=\sffamily\footnotesize] at (30:6.8) {\ang{60}};
  \draw[bkInk, line width=0.5pt] ($(B)+(-4,0)$) arc (180:105:4);
  \node[font=\sffamily\footnotesize] at ($(B)+(142:7)$) {\ang{75}};
  \draw[bkInk, line width=0.5pt] ($(C)+(240:5)$) arc (240:285:5);
  \node[font=\sffamily\footnotesize] at ($(C)+(262:8)$) {\ang{45}};
\end{tikzpicture}
\end{minipage}
\bkivhline{$b=\dfrac{c\sin B}{\sin C}=\dfrac{100\times\sin\ang{75}}{\sin\ang{45}}=\dfrac{100\times0.9659}{0.7071}\approx136.6$}
木までの距離は約\qty{137}{m}です（≈ は「ほぼ等しい」という記号。4-6、5-8）。川を渡らなくても、こちらの岸ではかった1つの長さと2つの角だけで、計算できました。

\bksubsection{4}{余弦定理}

\leavevmode\label{bk:4-7:yogen}\noindent\begin{minipage}[c]{0.58\linewidth}
池をはさんで、2本の木A、Bがあります。池の中は歩けないので、ABの長さを直接はかれません。そこで、両方が見える点Cに立ち、$\mathrm{CA}=\qty{80}{m}$、$\mathrm{CB}=\qty{60}{m}$、Cの角\ang{50}をはかりました。\par\vspace{2pt}
こんどは、わかっている角 $C$ の向かいの辺 $c$ がわかりません。正弦定理には「角と、その向かいの辺」の組が必要なので、このままでは使えません。新しい道具を作ります。
\end{minipage}\hfill
\begin{minipage}[c]{0.38\linewidth}\centering
\begin{tikzpicture}[x=1mm, y=1mm]
  \coordinate (C) at (0,0); \coordinate (B) at (36,0); \coordinate (A) at ({48*cos(50)},{48*sin(50)});
  \fill[bkBlue!20] ($(A)!0.5!(B)+(2,1)$) ellipse (9 and 6);
  \draw[bkInk, line width=0.8pt] (C) -- (B) (C) -- (A);
  \draw[bkRed, line width=1.1pt, dashed] (A) -- (B);
  \fill[bkLeaf] (A) circle (1.6) (B) circle (1.6);
  \fill[bkInk] (C) circle (0.8);
  \node[font=\sffamily\footnotesize, anchor=south] at ($(A)+(0,1.8)$) {A};
  \node[font=\sffamily\footnotesize, anchor=west] at ($(B)+(1.8,0)$) {B};
  \node[font=\sffamily\footnotesize, anchor=north] at ($(C)+(0,-1)$) {C};
  \node[font=\sffamily\footnotesize, anchor=north] at (18,-1) {\qty{60}{m}};
  \node[font=\sffamily\footnotesize, anchor=south east] at ($(C)!0.5!(A)+(-0.4,0)$) {\qty{80}{m}};
  \draw[bkInk, line width=0.5pt] (6,0) arc (0:50:6);
  \node[font=\sffamily\footnotesize] at (25:9) {\ang{50}};
  \node[font=\sffamily\footnotesize, text=bkBlue!70!black] at ($(A)!0.5!(B)+(2,1)$) {池};
\end{tikzpicture}
\end{minipage}

\bkivhstep{1}{座標と2点の距離で作る}
①の置き方を使います。Cを原点、Bを $(a,\ 0)$、Aを $(b\cos C,\ b\sin C)$ に置きます。$c$ はAとBの距離なので、2点の距離の式（4-4）で
\begin{align*}
c^2&=(b\cos C-a)^2+(b\sin C-0)^2 && \text{\small（2点の距離の2乗）}\\
&=b^2\cos^2C-2ab\cos C+a^2+b^2\sin^2C && \text{\small（展開する。3-8）}\\
&=a^2+b^2(\sin^2C+\cos^2C)-2ab\cos C && \text{\small（$b^2$ でくくる）}\\
&=a^2+b^2-2ab\cos C && \text{\small（$\sin^2C+\cos^2C=1$。4-6）}
\end{align*}
これを\textbf{余弦定理}（よげんていり）といいます。\bkidx{よげんていり}{余弦定理}
\bkivhline{\fbox{$c^2=a^2+b^2-2ab\cos C$}}
\textbf{2辺と、そのあいだの角}から、のこりの辺が求められます。ほかの辺についても、同じ形の式が成り立ちます。
\bkivhline{$a^2=b^2+c^2-2bc\cos A$\qquad $b^2=c^2+a^2-2ca\cos B$}
どの式でも、左の辺と、右の $\cos$ の角が向かい合っています。

\bkivhstep{2}{三平方の定理からの「ずれ」}
角 $C$ を\ang{90}にすると、$\cos\ang{90}=0$ なので $c^2=a^2+b^2$。三平方の定理にもどります。余弦定理は、三平方の定理を、直角でない三角形に広げた式なのです。$-2ab\cos C$ は、直角からのずれを直す部分です。

コンパスを開くように、$a=3$、$b=4$ のまま角 $C$ を変えると、$c$ は次の表のように変わります。$C$ が\ang{90}より小さいと $\cos C>0$ で、$c$ は5より短くなります。\ang{90}より大きいと $\cos C<0$ で、$-2ab\cos C$ が正になり、$c$ は5より長くなります。

\noindent\begin{minipage}[c]{0.46\linewidth}\centering
\begin{tikzpicture}[x=1mm, y=1mm]
  \def\s{7}
  \coordinate (C) at (0,0); \coordinate (B) at ({3*\s},0);
  \foreach \ang/\col in {60/bkBlue, 90/bkGreen, 120/bkOrange} {
    \draw[\col!80!black, line width=0.9pt] (C) -- ({4*\s*cos(\ang)},{4*\s*sin(\ang)});
    \draw[\col, line width=1.1pt, dashed] (B) -- ({4*\s*cos(\ang)},{4*\s*sin(\ang)});
    \fill[\col!80!black] ({4*\s*cos(\ang)},{4*\s*sin(\ang)}) circle (0.8);}
  \draw[bkInk, line width=1.2pt] (C) -- (B);
  \node[font=\sffamily\footnotesize, anchor=north] at ({1.5*\s},-1) {3};
  \node[font=\sffamily\footnotesize, text=bkBlue!80!black, anchor=west] at ({4*\s*cos(60)+1},{4*\s*sin(60)}) {\ang{60}};
  \node[font=\sffamily\footnotesize, text=bkGreen!50!black, anchor=south] at (0,{4*\s+0.6}) {\ang{90}};
  \node[font=\sffamily\footnotesize, text=bkOrange!80!black, anchor=east] at ({4*\s*cos(120)-1},{4*\s*sin(120)}) {\ang{120}};
  \node[font=\sffamily\footnotesize, anchor=north] at (0,-1) {C};
\end{tikzpicture}
\par\vspace{2pt}
{\footnotesize
\begin{tblr}{colspec={Q[c,m]Q[c,m]Q[c,m]}, row{1}={bg=bkTheme!12, font=\sffamily\footnotesize}, cells={font=\footnotesize}, hline{1,2,Z}={0.6pt,bkGray}, colsep=3pt, rowsep=0pt}
$C$ & $\cos C$ & $c=\sqrt{25-24\cos C}$ \\
\ang{0} & 1 & 1 \\
\ang{60} & 0.5 & $\sqrt{13}\approx3.61$ \\
\ang{90} & 0 & 5 \\
\ang{120} & $-0.5$ & $\sqrt{37}\approx6.08$ \\
\ang{180} & $-1$ & 7 \\
\end{tblr}}
\end{minipage}\hfill
\begin{minipage}[c]{0.5\linewidth}\raggedright
図は、長さ3の辺（黒）を横に置いたまま、長さ4の辺を\ang{60}・\ang{90}・\ang{120}に開いたところです。点線が、向かいの辺 $c$ です。\par\vspace{3pt}
開く角が大きいほど、点線は長くなります。\ang{90}のときが三平方の定理の $c=5$ で、それより閉じると短く、開くと長くなります。\par\vspace{3pt}
\ang{0}では、2辺が重なって $c=4-3=1$。\ang{180}では、一直線になって $c=3+4=7$。表のはしの値も、図で考えたとおりです。
\end{minipage}
\par\vspace{4pt}


\bkivhhead{池の答え}
$a=60$、$b=80$、$C=\ang{50}$ を入れます（$\cos\ang{50}\approx0.6428$）。
\begin{align*}
c^2&=60^2+80^2-2\times60\times80\times\cos\ang{50}\\
&=3600+6400-9600\times0.6428=\num{10000}-6170.9=3829.1
\end{align*}
$c=\sqrt{3829.1}\approx61.9$ なので、ABは約\qty{61.9}{m}です。

\bkivhstep{3}{3辺から角を求める}
余弦定理を $\cos C$ について解くと（2-14）、3辺の長さから角が求められます。$c^2=a^2+b^2-2ab\cos C$ で、$2ab\cos C$ を左へ、$c^2$ を右へ移すと $2ab\cos C=a^2+b^2-c^2$。両辺を $2ab$ でわって、
\bkivhline{\fbox{$\cos C=\dfrac{a^2+b^2-c^2}{2ab}$}}
たとえば、3辺が3、5、7の三角形で、いちばん長い辺7の向かいの角を $C$ とすると、
\bkivhline{$\cos C=\dfrac{3^2+5^2-7^2}{2\times3\times5}=\dfrac{9+25-49}{30}=\dfrac{-15}{30}=-\dfrac{1}{2}$}
$\cos C$ が負なので、$C$ は鈍角です。4-6⑧の表から $C=\ang{120}$ です。$\cos$ の値から、角は\ang{0}〜\ang{180}でただ1つに決まるので（4-6⑧）、安心して $\cos^{-1}$ が使えます。

\begin{bkfigure}
\begin{tikzpicture}[x=1mm, y=1mm]
  \foreach \xx/\aa/\cc/\lab/\cv in {0/72.54/bkBlue/{$c=5$}/{$\cos C=0.3$（正）\\鋭角 約\ang{72.5}}, 44/90/bkGreen!70!black/{$c=\sqrt{34}$}/{$\cos C=0$\\直角 \ang{90}}, 92/120/bkOrange/{$c=7$}/{$\cos C=-0.5$（負）\\鈍角 \ang{120}}} {
    \begin{scope}[shift={(\xx,0)}]
      \coordinate (C) at (0,0); \coordinate (B) at (15,0); \coordinate (A) at ({25*cos(\aa)},{25*sin(\aa)});
      \fill[\cc!12] (A) -- (B) -- (C) -- cycle;
      \draw[bkInk, line width=0.8pt] (C) -- (B) (C) -- (A);
      \draw[\cc, line width=1.4pt] (A) -- (B);
      \draw[bkInk, line width=0.5pt] (4,0) arc (0:\aa:4);
      \node[font=\sffamily\footnotesize, anchor=north] at (0,-1) {C};
      \node[font=\sffamily\footnotesize, anchor=north] at (7.5,-1) {3};
      \node[font=\sffamily\footnotesize, anchor=west, text=\cc!80!black] at ($(A)!0.5!(B)+(1.2,0.5)$) {\lab};
      \node[font=\sffamily\footnotesize, anchor=north, align=center] at (8,-6.5) {\cv};
    \end{scope}}
\end{tikzpicture}
\bkfigcaption{$a=3$、$b=5$（Cからのびる2辺）のまま、向かいの辺 $c$ を長くしていくと、$\cos C$ が正→0→負と変わり、角 $C$ が鋭角→直角→鈍角になる。}
\end{bkfigure}

\Needspace{26\baselineskip}
\begin{motto}[3]
\textbf{3辺から面積：ヘロンの公式と内接円}\par
3辺 $a$、$b$、$c$ がわかれば、余弦定理で $\cos C$、$\sin^2C+\cos^2C=1$ で $\sin C$ が求まり、$\dfrac{1}{2}ab\sin C$ で面積が出ます。この計算を文字のまま進めて1つの式にまとめたのが、\textbf{ヘロンの公式}です。$s=\dfrac{a+b+c}{2}$（周の長さの半分）として、
\bkivhline{$S=\sqrt{s(s-a)(s-b)(s-c)}$}
3辺が5、6、7なら、$s=9$ で、$S=\sqrt{9\times4\times3\times2}=\sqrt{216}\approx14.70$。$\cos C=\dfrac{25+36-49}{60}=0.2$、$\sin C=\sqrt{1-0.04}\approx0.9798$ から $\dfrac{1}{2}\times5\times6\times0.9798\approx14.70$ と求めても同じです。\par\vspace{2pt}
三角形の3つの辺にふれる円を\textbf{内接円}といいます。内接円の中心と3つの頂点を結ぶと、高さが半径 $r$ の三角形が3つできるので、$S=\dfrac{1}{2}r(a+b+c)$。5、6、7の三角形なら、$r=\dfrac{2S}{a+b+c}\approx\dfrac{29.39}{18}\approx1.63$ です。
\end{motto}

\bksubsection{5}{どの定理を使うか}

\leavevmode\label{bk:4-7:erabu}4-1で、「三角形が1つに決まる条件」（合同条件）を学びました。三角形が1つに決まるなら、のこりの辺や角も1つに決まるはずです。その求め方を、表にまとめます。

\begin{bkcenter}
\begin{bktable}{colspec={Q[l,m,wd=46mm] Q[l,m,wd=76mm]}, rowsep=2pt, row{1}={font=\sffamily\small}}
わかっているもの & 使う道具と順番 \\
3辺 & 余弦定理（$\cos C=\dfrac{a^2+b^2-c^2}{2ab}$）で角を求める。 \\
2辺とそのあいだの角 & 余弦定理（$c^2=a^2+b^2-2ab\cos C$）でのこりの辺。面積は $\dfrac{1}{2}ab\sin C$。 \\
1辺とその両はしの角 & 内角の和で3つめの角。次に正弦定理で辺。 \\
\end{bktable}
\bktabcaption{目じるしは「角と、その向かいの辺の組」がそろっているかどうか。そろっていれば正弦定理、なければ余弦定理。}
\end{bkcenter}

\begin{motto}[3]
\textbf{2辺と、そのあいだでない角がわかっているとき}\par
$a=5$、$c=8$、$A=\ang{30}$ のように、2辺と「そのあいだでない角」がわかっていても、合同条件ではないので、三角形は1つに決まるとはかぎりません。正弦定理で $\sin C=\dfrac{c\sin A}{a}=\dfrac{8\times0.5}{5}=0.8$。$\sin C=0.8$ になる角は、$C\approx\ang{53.1}$ と $C\approx\ang{126.9}$ の2つあります（4-6⑧）。どちらも内角の和が\ang{180}をこえないので、条件に合う三角形が2つできます。
\end{motto}

\bksubsection{6}{ロボットのうでを動かす：順運動学と逆運動学}

\leavevmode\label{bk:4-7:robot}4-0で、ロボットのうでについて、2つの問題を考えました。
\begin{itemize}[leftmargin=1.2\zw, itemsep=1pt, topsep=2pt, label={\color{\bkTT}●}]
\item \textbf{順運動学}：関節の角度から、手の位置を求める。
\item \textbf{逆運動学}：手を置きたい位置から、関節の角度を求める。
\end{itemize}
ここでは、肩とひじの2つの関節をもつ、平面の上で動くうで（2リンクのアーム）で、両方を計算します。

\noindent\begin{minipage}[c]{0.5\linewidth}\centering
\begin{tikzpicture}[x=1mm, y=1mm]
  \draw[bkGray, line width=0.6pt, -{Stealth[length=4pt]}] (-6,0) -- (50,0) node[anchor=north, font=\sffamily\footnotesize] {$x$};
  \draw[bkGray, line width=0.6pt, -{Stealth[length=4pt]}] (0,-4) -- (0,48) node[anchor=east, font=\sffamily\footnotesize] {$y$};
  \bkivharm{34}{25}{30}{60}{bkBlue!70}
  \draw[bkGray, line width=0.6pt, dashed] (E) -- ++(30:14);
  \draw[bkInk, line width=0.5pt] (8,0) arc (0:30:8);
  \node[font=\sffamily\footnotesize] at (15:11) {$\theta_1$};
  \draw[bkInk, line width=0.5pt] ($(E)+(30:7)$) arc (30:90:7);
  \node[font=\sffamily\footnotesize] at ($(E)+(60:9.6)$) {$\theta_2$};
  \node[font=\sffamily\footnotesize, anchor=north west] at ($(O)!0.5!(E)+(0.6,-0.4)$) {$L_1$};
  \node[font=\sffamily\footnotesize, anchor=east] at ($(E)!0.5!(P)+(-1.6,0)$) {$L_2$};
  \node[font=\sffamily\footnotesize, anchor=north west] at ($(O)+(1.6,-1.2)$) {O（肩）};
  \node[font=\sffamily\footnotesize, anchor=north west] at ($(E)+(1.2,-1)$) {E（ひじ）};
  \node[font=\sffamily\footnotesize, anchor=south] at ($(P)+(0,1.4)$) {P（手）};
\end{tikzpicture}
\end{minipage}\hfill
\begin{minipage}[c]{0.46\linewidth}
肩Oを原点にします。上のうでOEの長さを $L_1$、下のうでEPの長さを $L_2$ とします。
\begin{itemize}[leftmargin=1.2\zw, itemsep=1pt, topsep=2pt, label={\color{\bkTT}●}]
\item $\theta_1$：上のうでが、$x$ 軸からどれだけ回っているか。
\item $\theta_2$：下のうでが、上のうでをまっすぐのばした向き（点線）から、さらにどれだけ回っているか。
\end{itemize}
どちらも左回りを正とします。うでをまっすぐのばすと $\theta_2=\ang{0}$ です。
\end{minipage}

\bkivhstep{1}{順運動学：角度から手の位置}
ひじEは、Oから長さ $L_1$、角 $\theta_1$ の所なので（4-6⑧）、$\mathrm{E}(L_1\cos\theta_1,\ L_1\sin\theta_1)$ です。下のうでは、上のうでの向き $\theta_1$ から、さらに $\theta_2$ 回っているので、$x$ 軸から見ると $\theta_1+\theta_2$ の向きです（角はたし算で重なる）。手Pは、Eから長さ $L_2$、角 $\theta_1+\theta_2$ の向きに進んだ所なので、
\bkivhline{\fbox{$x=L_1\cos\theta_1+L_2\cos(\theta_1+\theta_2)$}\qquad \fbox{$y=L_1\sin\theta_1+L_2\sin(\theta_1+\theta_2)$}}
\bkidx{じゅんうんどうがく}{順運動学}たとえば $L_1=40$、$L_2=30$（cm）、$\theta_1=\ang{30}$、$\theta_2=\ang{60}$ なら、$\theta_1+\theta_2=\ang{90}$ なので、
\bkivhline{$x=40\cos\ang{30}+30\cos\ang{90}=34.6+0=34.6$\qquad $y=40\sin\ang{30}+30\sin\ang{90}=20+30=50$}
手は $(34.6,\ 50)$ にあります。下のうでがまっすぐ上（\ang{90}）を向いているので、$x$ はひじと同じ、$y$ はひじより30上、と図からも読めます。$\theta_1+\theta_2$ が\ang{90}をこえると $\cos(\theta_1+\theta_2)$ は負になり、手はひじより左へ出ます（\bkreidairef{7}）。この章では、$\theta_1+\theta_2$ が\ang{0}〜\ang{180}のときを考えます（それより大きい角は4-8）。

\bkivhstep{2}{逆運動学：手の位置から角度}
\bkidx{ぎゃくうんどうがく}{逆運動学}こんどは反対に、手を置きたい点 $\mathrm{P}(x,\ y)$ から、$\theta_1$ と $\theta_2$ を求めます。カギは、三角形OEPです。この三角形の3辺は、$L_1$、$L_2$、それにOPの長さ $d$ で、3つともわかります。3辺がわかれば、余弦定理で角が求められます（⑤の表の1行め）。

\begin{bkfigure}
\begin{tikzpicture}[x=1mm, y=1mm]
  \draw[bkGray, line width=0.6pt, -{Stealth[length=4pt]}] (-6,0) -- (52,0) node[anchor=north, font=\sffamily\footnotesize] {$x$};
  \draw[bkGray, line width=0.6pt, -{Stealth[length=4pt]}] (0,-4) -- (0,48) node[anchor=east, font=\sffamily\footnotesize] {$y$};
  % ひじを上にする形（点線）
  \draw[bkGray!80, line width=1.6pt, dashed, line cap=round] (0,0) -- ({32*cos(71.61)},{32*sin(71.61)}) -- (33.6,35.2);
  \bkivharm{32}{24}{21.05}{60}{bkBlue!70}
  \fill[bkOrange!12] (O) -- (E) -- (P) -- cycle;
  \draw[bkOrange, line width=0.9pt, dashed] (O) -- (P);
  \draw[bkBlue!70, line width=3.2pt, line cap=round] (O) -- (E) -- (P);
  \filldraw[fill=white, draw=bkInk, line width=0.7pt] (O) circle (1.6);
  \filldraw[fill=white, draw=bkInk, line width=0.7pt] (E) circle (1.4);
  \fill[bkInk] (P) circle (1.1);
  \draw[bkInk, line width=0.5pt] (10,0) arc (0:46.33:10);
  \node[font=\sffamily\footnotesize, anchor=west] at (8:10.6) {$\alpha$};
  \draw[bkRed, line width=0.6pt] (21.05:16) arc (21.05:46.33:16);
  \node[font=\sffamily\footnotesize, text=bkRed!80!black] at (33.7:18.6) {$\beta$};
  \node[font=\sffamily\footnotesize, text=bkOrange!80!black, anchor=south east] at ($(O)!0.5!(P)+(-0.5,0.5)$) {$d$};
  \node[font=\sffamily\footnotesize, anchor=north west] at ($(E)+(1,-1)$) {E};
  \node[font=\sffamily\footnotesize, anchor=south west] at ($(P)+(1,0.5)$) {$\mathrm{P}(x,\ y)$};
  \node[font=\sffamily\footnotesize, anchor=east] at ($(O)+(-4.5,1.5)$) {O};
  \node[font=\sffamily\footnotesize, text=bkGray, anchor=east, align=right] at (-1.5,28) {ひじを\\上にする形};
  % 説明
  \node[font=\sffamily\small, anchor=north west, align=left, text width=58mm] at (58,46) {%
  ① $d=\sqrt{x^2+y^2}$（4-4）\\[3pt]
  ② 三角形OEPで余弦定理。\\
  \quad ひじの所の角 $\angle\mathrm{OEP}$ を求める。\\
  \quad $\theta_2=\ang{180}-\angle\mathrm{OEP}$\\[3pt]
  ③ $\alpha$：OPと $x$ 軸の角。$\tan\alpha=\dfrac{y}{x}$\\
  \quad $\beta$：OEとOPの角（余弦定理）\\
  \quad $\theta_1=\alpha-\beta$};
\end{tikzpicture}
\bkfigcaption{三角形OEP（オレンジの面）の3辺 $L_1$、$L_2$、$d$ から、角を求める。点線は、もう1つの答え（ひじを上にする形）。}
\end{bkfigure}

\noindent なぜ $\theta_2=\ang{180}-\angle\mathrm{OEP}$ なのでしょう。うでをまっすぐのばすと、$\theta_2=\ang{0}$ で、ひじの所の角は\ang{180}です。下のうでを $\theta_2$ だけ曲げると、ひじの所の角は\ang{180}から $\theta_2$ だけ小さくなります。だから $\angle\mathrm{OEP}=\ang{180}-\theta_2$ です。

\bkivhhead{数で確かめる：手を $(40,\ 30)$ に置く（$L_1=40$、$L_2=30$）}
① $d=\sqrt{40^2+30^2}=\sqrt{2500}=50$\par
② ひじの所の角は、余弦定理で
\bkivhline{$\cos\angle\mathrm{OEP}=\dfrac{L_1^2+L_2^2-d^2}{2L_1L_2}=\dfrac{1600+900-2500}{2\times40\times30}=\dfrac{0}{2400}=0$}
$\angle\mathrm{OEP}=\ang{90}$ なので、$\theta_2=\ang{180}-\ang{90}=\ang{90}$\par
③ $\tan\alpha=\dfrac{30}{40}=0.75$ より $\alpha\approx\ang{36.87}$。三角形OEPで、Oの所の角 $\beta$ は
\bkivhline{$\cos\beta=\dfrac{L_1^2+d^2-L_2^2}{2L_1d}=\dfrac{1600+2500-900}{2\times40\times50}=\dfrac{3200}{4000}=0.8$\qquad $\beta\approx\ang{36.87}$}
だから $\theta_1=\alpha-\beta=\ang{0}$。上のうでは $x$ 軸の上にまっすぐのび（ひじは $(40,\ 0)$）、下のうでは\ang{90}曲がってまっすぐ上を向き、手は $(40,\ 30)$ に届きます。図をかいても、そのとおりです。

\bkivhhead{答えは2つあることがある}
点線のように、ひじを上にする形でも、同じ点に手が届きます。そのときは $\theta_1=\alpha+\beta$、$\theta_2$ は右回りの $-(\ang{180}-\angle\mathrm{OEP})$ です。人のうでも、同じ所に手を置くのに、ひじの位置を変えられますね。ロボットでは、ぶつかる物がないほうを選びます。また、$d$ が $L_1+L_2$ より大きい所には、うでをのばしきっても届かないので、答えはありません。

\bksubsection{7}{力の三角形：合力と、3つの力のつり合い}

\leavevmode\label{bk:4-7:chikara}4-6では、1つの力を長方形の2辺に分けました。反対に、2つの力を1つにまとめることもできます。

\bkivhstep{1}{合力の大きさ}
\noindent\begin{minipage}[c]{0.52\linewidth}
2人で、ひもをつけた荷物を引きます。1人は\qty{30}{N}、もう1人は\qty{40}{N}で、2本のひもの角は\ang{60}です。荷物は、2つの力を2辺とする平行四辺形の、対角線の向きに動きます。この対角線が、2つの力を合わせた1つの力で、\textbf{合力}（ごうりょく）といいます（4-6と同じく、実験で確かめられた法則。4-10でくわしく）。
\end{minipage}\hfill
\begin{minipage}[c]{0.44\linewidth}\centering
\begin{tikzpicture}[x=1mm, y=1mm]
  \coordinate (O) at (0,0); \coordinate (F1) at (30,0); \coordinate (F2) at ({28*cos(60)},{28*sin(60)});
  \coordinate (R) at ($(F1)+(F2)$);
  \fill[bkTheme!10] (O) -- (F1) -- (R) -- (F2) -- cycle;
  \draw[bkGray, line width=0.5pt, dashed] (F1) -- (R) -- (F2);
  \fill[bkYellow!35] (O) -- (F1) -- (R) -- cycle;
  \draw[bk arrow, bkBlue, line width=1.2pt] (O) -- (F1);
  \draw[bk arrow, bkRed, line width=1.2pt] (O) -- (F2);
  \draw[bk arrow, bkInk, line width=1.4pt] (O) -- (R);
  \draw[bkInk, line width=0.5pt] (5,0) arc (0:60:5);
  \node[font=\sffamily\footnotesize] at (28:8.6) {\ang{60}};
  \draw[bkInk, line width=0.5pt] ($(F1)+(60:5)$) arc (60:180:5);
  \node[font=\sffamily\footnotesize] at ($(F1)+(120:8)$) {\ang{120}};
  \node[font=\sffamily\footnotesize, text=bkBlue!80!black, anchor=north] at (15,-1) {\qty{30}{N}};
  \node[font=\sffamily\footnotesize, text=bkRed!80!black, anchor=east] at ($(O)!0.5!(F2)+(-0.8,0)$) {\qty{40}{N}};
  \node[font=\sffamily\footnotesize, anchor=west] at ($(F1)!0.5!(R)+(0.8,0)$) {40};
  \node[font=\sffamily\footnotesize, anchor=west] at ($(R)+(1,0)$) {合力};
\end{tikzpicture}
\end{minipage}

\par\vspace{4pt}
黄色の三角形を見ます。2辺は30と40（平行四辺形の向かい合う辺は等しい）で、そのあいだの角は\ang{120}です。平行四辺形のとなり合う角をたすと\ang{180}になるからです（平行線の角。4-1）。2辺とそのあいだの角がわかるので、余弦定理で
\begin{align*}
R^2&=30^2+40^2-2\times30\times40\times\cos\ang{120}\\
&=900+1600-2400\times(-0.5)=2500+1200=3700
\end{align*}
$R=\sqrt{3700}\approx60.8$ で、合力は約\qty{60.8}{N}です。$30+40=70$ より小さいのは、2人の引く向きが少しずれているからです。

\noindent 一般に、2つの力 $F_1$、$F_2$ のあいだの角を $\theta$ とすると、三角形の中の角は $\ang{180}-\theta$ です。$\cos(\ang{180}-\theta)=-\cos\theta$（4-6⑧）なので、余弦定理の $-2F_1F_2\cos(\ang{180}-\theta)$ は $+2F_1F_2\cos\theta$ になります。だから、
\bkivhline{\fbox{合力の大きさ\quad $R=\sqrt{F_1^2+F_2^2+2F_1F_2\cos\theta}$}}
（$\theta$ は2つの力のあいだの角）。
$\theta=\ang{0}$（同じ向き）なら $R=F_1+F_2$、$\theta=\ang{90}$ なら三平方の定理、$\theta=\ang{180}$（反対向き）なら $R$ は大きいほうから小さいほうをひいた差になります。

\bkivhstep{2}{3つの力のつり合いと正弦定理}
\noindent\begin{minipage}[c]{0.5\linewidth}
天井から、2本のロープでランプをつるします。ランプは止まっているので、ランプにはたらく3つの力（2本のロープが引く力と、重力 $W$）がつり合っています。\par\vspace{2pt}
つり合うには、2本のロープの力の合力が、重力と同じ大きさで、ちょうど真上を向いていなければなりません。合力の平行四辺形の半分の三角形（図\ref{bk:4-7:fig-lamp}の黄色）を見ると、3辺が、2本のロープの力と $W$ の三角形になります。
\end{minipage}\hfill
\begin{minipage}[c]{0.46\linewidth}\centering
\begin{tikzpicture}[x=1mm, y=1mm]
  \def\pa{30}\def\pb{45}
  \coordinate (K) at (0,0);
  \coordinate (U1) at ({-34*sin(\pa)},{34*cos(\pa)});
  \coordinate (U2) at ({34*sin(\pb)*0.9},{34*cos(\pb)*0.9+3.5});
  \fill[bkGray!40] (-24,{34*cos(\pa)}) rectangle (26,{34*cos(\pa)+2});
  \draw[bkInk!70, line width=0.8pt] (K) -- (U1) (K) -- (U2);
  % 力の矢印（長さは力の大きさに比例： 1 mm = 3.5 N）
  \coordinate (T1) at ({-71.7/3.5*sin(\pa)},{71.7/3.5*cos(\pa)});
  \coordinate (T2) at ({50.7/3.5*sin(\pb)},{50.7/3.5*cos(\pb)});
  \coordinate (Up) at (0,{98/3.5});
  \fill[bkYellow!40] (K) -- (T1) -- (Up) -- cycle;
  \draw[bkGray, line width=0.5pt, dashed] (T1) -- (Up) -- (T2);
  \draw[bk arrow, bkBlue, line width=1.2pt] (K) -- (T1);
  \draw[bk arrow, bkRed, line width=1.2pt] (K) -- (T2);
  \draw[bk arrow, bkInk, line width=1.2pt] (K) -- (0,{-98/3.5});
  \draw[bkGray, line width=0.6pt, densely dotted] (K) -- (Up);
  \fill[bkYellow!80!black] (-3,-4) rectangle (3,-1) ;
  \fill[bkInk] (K) circle (0.8);
  \node[font=\sffamily\footnotesize, text=bkBlue!80!black, anchor=east] at ($(T1)+(-0.6,0)$) {$T_{\mathrm{L}}$};
  \node[font=\sffamily\footnotesize, text=bkRed!80!black, anchor=west] at ($(T2)+(0.6,0)$) {$T_{\mathrm{R}}$};
  \node[font=\sffamily\footnotesize, anchor=west] at (1,{-98/3.5+2}) {$W$};
  \draw[bkInk, line width=0.5pt] (0,6) arc (90:{90+\pa}:6);
  \node[font=\sffamily\footnotesize, anchor=south] at ({90+\pa/2}:6.6) {$\theta_{\mathrm{L}}$};
  \draw[bkInk, line width=0.5pt] (0,9) arc (90:{90-\pb}:9);
  \node[font=\sffamily\footnotesize, anchor=south west] at ({90-\pb/2}:9) {$\theta_{\mathrm{R}}$};
\end{tikzpicture}
\bkfigcaption{ランプにはたらく3つの力。黄色が力の三角形。}\label{bk:4-7:fig-lamp}
\end{minipage}

\par\vspace{4pt}
\noindent 左のロープが真上からかたむいた角を $\theta_{\mathrm{L}}$、右のロープの角を $\theta_{\mathrm{R}}$ とします（L は左 left、R は右 right の頭文字）。ロープが引く力も、$T_{\mathrm{L}}$、$T_{\mathrm{R}}$ と書きます。黄色の三角形の角は、$W$ の向かいが $\ang{180}-(\theta_{\mathrm{L}}+\theta_{\mathrm{R}})$、$T_{\mathrm{L}}$ の向かいが $\theta_{\mathrm{R}}$、$T_{\mathrm{R}}$ の向かいが $\theta_{\mathrm{L}}$ です（平行四辺形の向かい合う辺は平行なので、平行線の錯角が等しい。4-1）。正弦定理を使うと、$\sin(\ang{180}-(\theta_{\mathrm{L}}+\theta_{\mathrm{R}}))=\sin(\theta_{\mathrm{L}}+\theta_{\mathrm{R}})$ なので、
\bkivhline{\fbox{$\dfrac{T_{\mathrm{L}}}{\sin\theta_{\mathrm{R}}}=\dfrac{T_{\mathrm{R}}}{\sin\theta_{\mathrm{L}}}=\dfrac{W}{\sin(\theta_{\mathrm{L}}+\theta_{\mathrm{R}})}$}}
となります。ロープを左右に大きく開くと、$\theta_{\mathrm{L}}+\theta_{\mathrm{R}}$ が\ang{180}に近づき、$\sin(\theta_{\mathrm{L}}+\theta_{\mathrm{R}})$ が0に近づくので、$T_{\mathrm{L}}$、$T_{\mathrm{R}}$ はとても大きくなります。洗たく物のロープを、ぴんとまっすぐには張れないのは、このためです。

\bkivhhead{力の三角形のまとめ}
力の矢印をつないで三角形を作ると、この章の2つの定理が、そのまま力の計算に使えます。
\begin{itemize}[leftmargin=1.2\zw, itemsep=1pt, topsep=2pt, label={\color{\bkTT}●}]
\item 2つの力から合力を求める：2辺とそのあいだの角がわかるので、余弦定理。
\item 3つの力のつり合いで、力の大きさを求める：角と向かいの辺の組があるので、正弦定理。
\end{itemize}
4-10・4-11では、力を矢印（ベクトル）として、たてと横の成分に分けて計算する方法を学びます。成分を使うと、三角形をかかなくても、たし算だけで合力が求められます。どちらの方法でも、同じ答えになります。

\begin{column}{4-0の\bkhatenaref に答える}
4-0では「うでの長さと関節の角度がわかれば、手の位置は決まる？　反対に、手を置きたい場所から角度は決まる？」と問いかけました。答えは次のとおりです。
\begin{itemize}[leftmargin=1.2\zw, itemsep=1pt, topsep=2pt, label={\color{\bkTT}●}]
\item 角度から位置（順運動学）は、三角比のたし算で、いつも1つに決まる。
\item 位置から角度（逆運動学）は、余弦定理で求める。答えは、ふつう2つ（ひじを上か下か）。届かない所では答えがない。
\end{itemize}
みなさんのうでも、肩とひじの2つの関節をもつ2リンクのうでとみなせます。机の上の1点に指先を置いたまま、ひじを上げ下げしてみましょう。ひじの位置が2通りあることが、体で確かめられます。工場の溶接ロボットは、1秒に何百回もこの計算をくり返して、手の先を決められた道すじにそって動かしています。関節が6つあるロボットの計算は、行列（第7巻）を使って整理します。
\end{column}

\begin{taisetsu}
三角形ABC：角 $A$ の向かいの辺が $a$。Cを原点、Bを $(a,\ 0)$ に置くと、$\mathrm{A}(b\cos C,\ b\sin C)$。\par\vspace{3pt}
面積：$S=\dfrac{1}{2}ab\sin C$（2辺とそのあいだの角）\par\vspace{3pt}
正弦定理：$\dfrac{a}{\sin A}=\dfrac{b}{\sin B}=\dfrac{c}{\sin C}=2R$（$R$ は外接円の半径）\par\vspace{3pt}
余弦定理：$c^2=a^2+b^2-2ab\cos C$\qquad $\cos C=\dfrac{a^2+b^2-c^2}{2ab}$\par\vspace{3pt}
ロボットのうで（順運動学）：\par\hspace*{2\zw}$x=L_1\cos\theta_1+L_2\cos(\theta_1+\theta_2)$\quad $y=L_1\sin\theta_1+L_2\sin(\theta_1+\theta_2)$
\end{taisetsu}

\begin{kotoba}
\small
\bkword{三角形を解く}{三角形のわかっている辺や角から、のこりの辺や角を求めること。}
\bkword{外接円}{三角形の3つの頂点を通る円。中心は外心。}
\bkword{正弦定理}{辺と、向かいの角の $\sin$ の比が、どの辺でも同じ（外接円の直径 $2R$）という定理。}
\bkword{余弦定理}{2辺とそのあいだの角から、のこりの辺を求める定理。$C=\ang{90}$ で三平方の定理になる。}
\bkword{順運動学・逆運動学}{関節の角度から手の位置を求めること・手の位置から関節の角度を求めること。}
\bkword{合力}{2つ以上の力を合わせて、同じ働きをする1つの力にしたもの。}
\end{kotoba}

%% ---------------- 例題 ----------------
\bksection{reidai}{例題でたしかめよう}

\begin{reidai}[2辺とあいだの角から面積]
三角形ABCの面積を求めましょう。\quad（1）$a=6$、$b=5$、$C=\ang{30}$\quad（2）$a=8$、$b=5$、$C=\ang{135}$\par
\kangaekata 2辺とそのあいだの角なので、$S=\dfrac{1}{2}ab\sin C$。（2）は鈍角なので、$\sin\ang{135}=\sin(\ang{180}-\ang{45})=\sin\ang{45}$ を使います。\par
\shiki（1）$S=\dfrac{1}{2}\times6\times5\times\sin\ang{30}=\dfrac{1}{2}\times6\times5\times\dfrac{1}{2}=\dfrac{30}{4}=7.5$\par
\hspace*{2.6\zw}（2）$S=\dfrac{1}{2}\times8\times5\times\sin\ang{135}=20\times\dfrac{\sqrt2}{2}=10\sqrt2\approx14.14$\par
\kotae（1）7.5\quad（2）$10\sqrt2$（約14.1）
\end{reidai}

\begin{jibun}
三角形ABCの面積を求めましょう。\par
（1）$a=4$、$b=7$、$C=\ang{150}$\par
（2）$a=6$、$b=4$、$C=\ang{60}$
\bkjans{（1）\nobreak $\dfrac{1}{2}\times4\times7\times\dfrac{1}{2}=7$。（2）\nobreak $6\sqrt3$（約10.4）。}
\end{jibun}
\textbf{ポイント}\quad 面積の式に入れる角は、2辺の「あいだ」の角です。鈍角でも $\sin$ は正なので、面積はいつも正になります。\bkreidairef{1}の（2）の三角形は、あいだの角を\ang{45}にした三角形と、高さも面積も同じです（②の\ang{60}と\ang{120}の図と同じ考え方）。

\begin{reidai}[正弦定理で辺と外接円の半径]
三角形ABCで、$a=10$、$A=\ang{45}$、$B=\ang{60}$ です。$b$ と、外接円の半径 $R$ を求めましょう。\par
\kangaekata 角 $A$ と向かいの辺 $a$ の組がそろっているので、正弦定理。$\dfrac{a}{\sin A}=\dfrac{b}{\sin B}$ の両辺に $\sin B$ をかけて $b$ を出します。\par
\shiki $b=\dfrac{a\sin B}{\sin A}=\dfrac{10\times\frac{\sqrt3}{2}}{\frac{\sqrt2}{2}}=\dfrac{10\sqrt3}{\sqrt2}=\dfrac{10\sqrt3\times\sqrt2}{\sqrt2\times\sqrt2}=\dfrac{10\sqrt6}{2}=5\sqrt6$\par
\hspace*{2.6\zw}$2R=\dfrac{a}{\sin A}=\dfrac{10}{\frac{\sqrt2}{2}}=\dfrac{20}{\sqrt2}=10\sqrt2$\quad より\quad $R=5\sqrt2$\par
\kotae $b=5\sqrt6$（約12.2）、$R=5\sqrt2$（約7.07）
\end{reidai}

\begin{jibun}
三角形ABCで、$a=8$、$A=\ang{30}$、$B=\ang{45}$ です。$b$ と、外接円の半径 $R$ を求めましょう。
\bkjans{$b=\dfrac{8\sin\ang{45}}{\sin\ang{30}}=8\sqrt2$（約11.3）、$R=8$。}
\end{jibun}
\textbf{ポイント}\quad 答えを確かめるには、大きい角の向かいに長い辺があるかを見ます。\bkreidairef{2}では $B=\ang{60}$ が $A=\ang{45}$ より大きいので、$b$ は $a=10$ より長いはずです。$b\approx12.2$ で、合っています。また、$\dfrac{b}{\sin B}=\dfrac{12.25}{0.866}\approx14.14$ と $\dfrac{a}{\sin A}=\dfrac{10}{0.7071}\approx14.14$ が等しいことも確かめられます。この値が外接円の直径 $2R=10\sqrt2\approx14.14$ です。

\begin{reidai}[1辺と両はしの角から、遠くの距離]
\noindent\begin{minipage}[t]{0.62\linewidth}
ABの長さは\qty{80}{m}です。Aから見た角は\ang{45}、Bから見た角は\ang{105}でした。ACとBCの長さを、小数第1位まで求めましょう。\par
\kangaekata 1辺とその両はしの角なので、まず内角の和で $C$ を出し、次に正弦定理（⑤の表の3行め）。\ang{105}は鈍角ですが、$\sin$ は電卓でそのまま計算できます。
\end{minipage}\hfill
\begin{minipage}[t]{0.34\linewidth}\vspace{0pt}\centering
\begin{tikzpicture}[x=1mm, y=1mm]
  \coordinate (A) at (0,0); \coordinate (B) at (22,0);
  \coordinate (C) at ({42.5*cos(45)},{42.5*sin(45)});
  \fill[bkTheme!12] (A) -- (B) -- (C) -- cycle;
  \draw[bkInk, line width=0.8pt] (A) -- (B) -- (C) -- cycle;
  \node[font=\sffamily\footnotesize, anchor=north] at (A) {A};
  \node[font=\sffamily\footnotesize, anchor=north] at (B) {B};
  \node[font=\sffamily\footnotesize, anchor=south] at ($(C)+(0,0.5)$) {C};
  \node[font=\sffamily\footnotesize, anchor=north] at (11,-1) {80};
  \draw[bkInk, line width=0.5pt] (5,0) arc (0:45:5);
  \node[font=\sffamily\footnotesize] at (22.5:8) {\ang{45}};
  \draw[bkInk, line width=0.5pt] ($(B)+(4,0)$) arc (0:75:4);
  \node[font=\sffamily\footnotesize, anchor=west] at ($(B)+(30:4.6)$) {\ang{105}};
\end{tikzpicture}
\end{minipage}\par
\shiki $C=\ang{180}-\ang{45}-\ang{105}=\ang{30}$\par
\hspace*{2.6\zw}$\mathrm{AC}=b=\dfrac{c\sin B}{\sin C}=\dfrac{80\times\sin\ang{105}}{\sin\ang{30}}=\dfrac{80\times0.9659}{0.5}\approx154.5$\par
\hspace*{2.6\zw}$\mathrm{BC}=a=\dfrac{c\sin A}{\sin C}=\dfrac{80\times\sin\ang{45}}{\sin\ang{30}}=\dfrac{80\times0.7071}{0.5}\approx113.1$\par
\kotae AC約\qty{154.5}{m}、BC約\qty{113.1}{m}
\end{reidai}
図の角Bは、ABをBの先へのばした線から見ると\ang{75}です。ABとBCのあいだの角は $\ang{180}-\ang{75}=\ang{105}$ です。

\begin{jibun}
ABの長さが\qty{50}{m}、Aから見た角が\ang{70}、Bから見た角が\ang{50}の三角形ABCで、ACとBCの長さを小数第1位まで求めましょう。
\bkjans{$C=\ang{60}$。AC $=\dfrac{50\sin\ang{50}}{\sin\ang{60}}\approx\qty{44.2}{m}$、BC $=\dfrac{50\sin\ang{70}}{\sin\ang{60}}\approx\qty{54.3}{m}$。}
\end{jibun}
\textbf{ポイント}\quad 1辺と両はしの角がわかっているときは、内角の和で3つめの角を出してから、正弦定理を使う、の2段です。答えを出したら、いちばん大きい角（\ang{105}）の向かいの辺（AC）がいちばん長いか確かめましょう。154.5、113.1、80の順で、合っています。

\begin{reidai}[余弦定理で辺]
三角形ABCで、$c$ を求めましょう。\quad（1）$a=5$、$b=8$、$C=\ang{60}$\quad（2）$a=3$、$b=5$、$C=\ang{120}$\par
\kangaekata 2辺とそのあいだの角がわかっているので、余弦定理を使います。（2）は $\cos\ang{120}=-\dfrac{1}{2}$ なので、ひき算がたし算に変わります。\par
\shiki（1）$c^2=5^2+8^2-2\times5\times8\times\dfrac{1}{2}=25+64-40=49$\par
\hspace*{4.6\zw}$c>0$ なので $c=7$\par
\hspace*{2.6\zw}（2）$c^2=3^2+5^2-2\times3\times5\times\Bigl(-\dfrac{1}{2}\Bigr)=9+25+15=49$\par
\hspace*{4.6\zw}$c=7$\par
\kotae（1）7\quad（2）7
\end{reidai}
（1）と（2）は、ちがう三角形なのに、どちらも $c=7$ になりました。$-2ab\cos C$ の部分が、（1）は $-40$、（2）は $+15$ で、ちょうど49になるように選んだ数です。

\begin{jibun}
三角形ABCで、$c$ を求めましょう。\par
（1）$a=5$、$b=3$、$C=\ang{60}$\par
（2）$a=4$、$b=6$、$C=\ang{120}$
\bkjans{（1）\nobreak $c^2=25+9-15=19$、$c=\sqrt{19}$（約4.36）。（2）\nobreak $c^2=16+36+24=76$、$c=\sqrt{76}=2\sqrt{19}$（約8.72）。}
\end{jibun}
\textbf{ポイント}\quad $\cos C$ の符号に気をつけましょう。あいだの角が鋭角なら $-2ab\cos C$ はひき算、鈍角ならたし算になります。答えの $c$ は、鋭角なら三平方の定理の $\sqrt{a^2+b^2}$ より短く、鈍角なら長くなります。（1）は $\sqrt{89}\approx9.4$ より短い7、（2）は $\sqrt{34}\approx5.8$ より長い7です。

\begin{reidai}[3辺から角]
3辺が $a=7$、$b=5$、$c=3$ の三角形ABCで、いちばん大きい角を求めましょう。\par
\kangaekata いちばん長い辺の向かいの角が、いちばん大きい角です（正弦定理）。$a=7$ がいちばん長いので、角 $A$ を求めます。$A$ の向かいの辺は $a$ なので、$\cos A=\dfrac{b^2+c^2-a^2}{2bc}$ の形を使います。\par
\shiki $\cos A=\dfrac{5^2+3^2-7^2}{2\times5\times3}=\dfrac{25+9-49}{30}=\dfrac{-15}{30}=-\dfrac{1}{2}$\par
\hspace*{2.6\zw}$\cos A$ が負なので鈍角。$\cos\ang{120}=-\dfrac{1}{2}$ より $A=\ang{120}$\par
\kotae $A=\ang{120}$
\end{reidai}

\begin{jibun}
3辺が $a=7$、$b=8$、$c=3$ の三角形ABCで、角 $A$ を求めましょう。
\bkjans{$\cos A=\dfrac{64+9-49}{48}=\dfrac{1}{2}$、$A=\ang{60}$。}
\end{jibun}
\textbf{ポイント}\quad 3辺から角を求めるときは、$\cos$ の式を使えば、角が1つに決まります（$\sin$ だと、鋭角と鈍角の2つが出ることがある）。いちばん長い辺の向かいの角の $\cos$ が、負なら鈍角三角形、0なら直角三角形、正なら鋭角三角形です。\bkreidairef{5}で、のこりの角も求めてみましょう。$\cos B=\dfrac{a^2+c^2-b^2}{2ac}=\dfrac{49+9-25}{42}=\dfrac{33}{42}\approx0.786$ より $B\approx\ang{38.2}$、内角の和から $C\approx\ang{180}-\ang{120}-\ang{38.2}=\ang{21.8}$ です。

\begin{reidai}[池をはさんだ2点の距離]
池をはさんだ2点A、Bの距離を知るために、点Cで $\mathrm{CA}=\qty{80}{m}$、$\mathrm{CB}=\qty{60}{m}$、Cの角\ang{50}をはかりました。ABを小数第1位まで求めましょう（④の場面）。\par
\kangaekata 2辺とそのあいだの角なので、余弦定理。$\cos\ang{50}$ は電卓で求めます。\par
\shiki $c^2=60^2+80^2-2\times60\times80\times\cos\ang{50}$\par
\hspace*{4.6\zw}$=3600+6400-9600\times0.64279$\par
\hspace*{4.6\zw}$=\num{10000}-6170.8=3829.2$\par
\hspace*{2.6\zw}$c=\sqrt{3829.2}\approx61.88$\par
\kotae 約\qty{61.9}{m}
\end{reidai}

\begin{jibun}
点Cから、$\mathrm{CA}=\qty{120}{m}$、$\mathrm{CB}=\qty{90}{m}$、Cの角が\ang{70}でした。ABを小数第1位まで求めましょう。
\bkjans{$c^2=\num{22500}-\num{21600}\cos\ang{70}\approx\num{15112.4}$、$c\approx\qty{122.9}{m}$。}
\end{jibun}
\textbf{ポイント}\quad 実際の測量では、はかった角の小さなずれも考えます。\bkreidairef{6}で、Cの角が\ang{51}なら約\qty{62.9}{m}、\ang{49}なら約\qty{60.8}{m}です。角が\ang{1}ずれると、答えは約\qty{1}{m}変わります。

\begin{reidai}[ロボットのうで：順運動学]
$L_1=40$、$L_2=30$（cm）の2リンクのうでで、$\theta_1=\ang{30}$、$\theta_2=\ang{90}$ のとき、ひじEと手Pの座標を、小数第1位まで求めましょう。\par
\kangaekata ひじは $(L_1\cos\theta_1,\ L_1\sin\theta_1)$。下のうでは $\theta_1+\theta_2=\ang{120}$ の向きなので、$\cos\ang{120}$ は負で、手はひじより左へ出ます。\par
\shiki ひじ：$x=40\cos\ang{30}=40\times0.8660=34.64$\par
\hspace*{5.4\zw}$y=40\sin\ang{30}=40\times0.5=20$\par
\hspace*{2.6\zw}手：$x=34.64+30\cos\ang{120}=34.64+30\times(-0.5)=19.64$\par
\hspace*{4.6\zw}$y=20+30\sin\ang{120}=20+30\times0.8660=45.98$\par
\kotae ひじ $(34.6,\ 20.0)$、手 $(19.6,\ 46.0)$
\end{reidai}

\begin{jibun}
$L_1=40$、$L_2=30$（cm）の2リンクのうでで、$\theta_1=\ang{60}$、$\theta_2=\ang{60}$ のとき、ひじと手の座標を小数第1位まで求めましょう。
\bkjans{ひじ $(20.0,\ 34.6)$。$\theta_1+\theta_2=\ang{120}$ より、手は $x=20-15=5.0$、$y=34.64+25.98\approx60.6$。}
\end{jibun}
\textbf{ポイント}\quad $\theta_1+\theta_2$ が\ang{90}をこえると、下のうでは左上を向き、$\cos(\theta_1+\theta_2)$ が負になります。手の $x$ 座標がひじより小さくなるのは、そのためです。図をかいて、手がひじより左にあるか確かめましょう。

\begin{reidai}[ロボットのうで：逆運動学]
$L_1=40$、$L_2=30$（cm）の2リンクのうでで、手を $\mathrm{P}(42,\ 44)$ に置きます。ひじを下にする形の $\theta_2$ と $\theta_1$ を、小数第1位まで求めましょう。\par
\kangaekata ⑥の①②③の順。$d$ を出し、三角形OEPで余弦定理を2回使います。\par
\shiki ① $d^2=42^2+44^2=1764+1936=3700$、$d=\sqrt{3700}\approx60.83$\par
\hspace*{2.6\zw}② $\cos\angle\mathrm{OEP}=\dfrac{1600+900-3700}{2\times40\times30}=\dfrac{-1200}{2400}=-0.5$\par
\hspace*{4.6\zw}$\angle\mathrm{OEP}=\ang{120}$、$\theta_2=\ang{180}-\ang{120}=\ang{60}$\par
\hspace*{2.6\zw}③ $\alpha=\tan^{-1}\dfrac{44}{42}\approx\ang{46.33}$\par
\hspace*{4.6\zw}$\cos\beta=\dfrac{1600+3700-900}{2\times40\times60.83}=\dfrac{4400}{4866.4}\approx0.9042$\par
\hspace*{4.6\zw}$\beta\approx\ang{25.28}$\par
\hspace*{4.6\zw}$\theta_1=\alpha-\beta\approx\ang{46.33}-\ang{25.28}=\ang{21.05}$\par
\kotae $\theta_1\approx\ang{21.1}$、$\theta_2=\ang{60}$
\end{reidai}
確かめ：順運動学の式に $\theta_1=\ang{21.05}$、$\theta_1+\theta_2=\ang{81.05}$ を入れると、
\bkivhline{\begin{tabular}{@{}l@{}}$x=40\cos\ang{21.05}+30\cos\ang{81.05}\approx37.33+4.67=42.0$\\[2pt] $y=40\sin\ang{21.05}+30\sin\ang{81.05}\approx14.37+29.63=44.0$\end{tabular}}
ねらった点に届いています。

\begin{jibun}
$L_1=40$、$L_2=30$（cm）の2リンクのうでで、手を $(20,\ 30)$ に置きます。ひじを下にする形の $\theta_2$ と $\theta_1$ を小数第1位まで求めましょう。
\bkjans{$d^2=1300$。$\cos\angle\mathrm{OEP}=0.5$ より $\theta_2=\ang{120}$。$\alpha\approx\ang{56.31}$、$\cos\beta\approx0.6934$ より $\beta\approx\ang{46.10}$。$\theta_1\approx\ang{10.2}$。}
\end{jibun}
\textbf{ポイント}\quad 逆運動学の答えは、順運動学の式に入れて、ねらった点に届くかで確かめられます。また、$d$ が $L_1+L_2$（ここでは70）より大きい点や、$L_1-L_2$（ここでは10）より小さい点には、手が届きません。余弦定理で $\cos$ の値が1より大きく（または $-1$ より小さく）なってしまうので、すぐにわかります。

\begin{reidai}[合力の大きさ]
1つの荷物を、\qty{30}{N}と\qty{40}{N}の力で引きます。2つの力のあいだの角が\ang{60}のとき、合力の大きさを小数第1位まで求めましょう。\par
\kangaekata 平行四辺形の半分の三角形で考えます。三角形の中の角は $\ang{180}-\ang{60}=\ang{120}$。余弦定理を使います。\par
\shiki $R^2=30^2+40^2-2\times30\times40\times\cos\ang{120}$\par
\hspace*{4.6\zw}$=900+1600-2400\times(-0.5)=3700$\par
\hspace*{2.6\zw}$R=\sqrt{3700}\approx60.83$\par
\kotae 約\qty{60.8}{N}
\end{reidai}

\begin{jibun}
1つの荷物を、\qty{50}{N}と\qty{50}{N}の力で引きます。2つの力のあいだの角が\ang{120}のとき、合力の大きさを求めましょう。
\bkjans{$R^2=2500+2500-5000\cos\ang{60}=2500$、$R=\qty{50}{N}$。}
\end{jibun}
\textbf{ポイント}\quad 2つの力のあいだの角が大きいほど、合力は小さくなります。\qty{50}{N}と\qty{50}{N}でも、あいだの角が\ang{0}なら\qty{100}{N}、\ang{90}なら約\qty{70.7}{N}、\ang{120}なら\qty{50}{N}、\ang{180}なら0です。

\begin{reidai}[2本のロープでつるす]
質量\qty{10}{kg}のランプ（重力 $W=\qty{98}{N}$）を、2本のロープでつるします。ロープが真上からかたむいた角は、$\theta_{\mathrm{L}}=\ang{30}$、$\theta_{\mathrm{R}}=\ang{45}$ です。2本のロープが引く力 $T_{\mathrm{L}}$、$T_{\mathrm{R}}$ を、小数第1位まで求めましょう。\par
\kangaekata ⑦の②の正弦定理の式 $\dfrac{T_{\mathrm{L}}}{\sin\theta_{\mathrm{R}}}=\dfrac{T_{\mathrm{R}}}{\sin\theta_{\mathrm{L}}}=\dfrac{W}{\sin(\theta_{\mathrm{L}}+\theta_{\mathrm{R}})}$ を使います。$\theta_{\mathrm{L}}+\theta_{\mathrm{R}}=\ang{75}$ です。\par
\shiki $\dfrac{W}{\sin\ang{75}}=\dfrac{98}{0.96593}\approx101.46$\par
\hspace*{2.6\zw}$T_{\mathrm{L}}=101.46\times\sin\ang{45}=101.46\times0.70711\approx71.74$\par
\hspace*{2.6\zw}$T_{\mathrm{R}}=101.46\times\sin\ang{30}=101.46\times0.5=50.73$\par
\kotae $T_{\mathrm{L}}$ 約\qty{71.7}{N}、$T_{\mathrm{R}}$ 約\qty{50.7}{N}
\end{reidai}
真上に近いロープ（$\theta_{\mathrm{L}}=\ang{30}$）のほうが、大きな力で引いています。確かめ：横向きの力は、$71.74\times\sin\ang{30}\approx35.9$ と $50.73\times\sin\ang{45}\approx35.9$ で、左右の引き合いがつり合っています。

\begin{jibun}
重力 $W=\qty{60}{N}$ の物を2本のロープでつるします。$\theta_{\mathrm{L}}=\theta_{\mathrm{R}}=\ang{60}$ のとき、$T_{\mathrm{L}}$、$T_{\mathrm{R}}$ を求めましょう。
\bkjans{$T_{\mathrm{L}}=T_{\mathrm{R}}=\dfrac{60\sin\ang{60}}{\sin\ang{120}}=\qty{60}{N}$（重力と同じ）。}
\end{jibun}
\textbf{ポイント}\quad 左右のロープを同じ角 $\theta$ でかたむけて、重力\qty{60}{N}の物をつるすとき、ロープ1本が引く力は、$\theta=\ang{30}$ なら約\qty{34.6}{N}、\ang{45}なら約\qty{42.4}{N}、\ang{60}なら\qty{60}{N}です（どれも正弦定理で計算）。ロープを開くほど、ロープには大きな力がかかります。

%% ---------------- まちがえやすいところ ----------------
\bksection{tsumazuki}{まちがえやすいところをチェック}

\begin{tsumazuki}{余弦定理で、あいだでない角を使う}
\bkng{$a=5$、$b=8$、$A=\ang{60}$ なので $c^2=5^2+8^2-2\times5\times8\cos\ang{60}$}
\bkok{$c^2=a^2+b^2-2ab\cos C$ の角は、$a$ と $b$ のあいだの角 $C$}
余弦定理の $\cos$ の角は、左の辺 $c$ の向かいの角で、右の2辺 $a$、$b$ のあいだの角です。図をかいて、角がどの2辺にはさまれているか確かめましょう。
\end{tsumazuki}

\begin{tsumazuki}{鈍角の $\cos$ を正にする}
\bkng{$\cos\ang{120}=0.5$ なので $c^2=9+25-15=19$}
\bkok{$\cos\ang{120}=-0.5$ なので $c^2=9+25+15=49$}
鈍角の $\cos$ は負です（4-6⑧）。鈍角の三角形では、向かいの辺は三平方の定理のときより長くなるはずです。$19<34$ では短くなってしまい、合いません。
\end{tsumazuki}

\begin{tsumazuki}{辺の比を、角の比にする}
\bkng{角が\ang{30}・\ang{60}・\ang{90}なので、辺の比は $1:2:3$}
\bkok{辺の比は $\sin\ang{30}:\sin\ang{60}:\sin\ang{90}=1:\sqrt3:2$}
正弦定理で、辺は角の $\sin$ に比例します。角そのものには比例しません。
\end{tsumazuki}

\begin{tsumazuki}{合力で、力のあいだの角をそのまま使う}
\bkng{\ang{60}で引く\qty{30}{N}と\qty{40}{N}の合力は $R^2=900+1600-2400\cos\ang{60}$}
\bkok{三角形の中の角 \ang{120} を使う。$R^2=900+1600-2400\cos\ang{120}$}
×の式では $R=\sqrt{1300}\approx36$ となり、40より小さくなります。同じような向きに引いて、片方の力より小さくなるのはおかしいですね。
\end{tsumazuki}

%% ---------------- 物理・工学 ----------------
\begin{butsuri}
\par\addvspace{3pt}\textbf{屋根のトラス：部材の長さを決める}\par\noindent
\begin{minipage}[c]{0.6\linewidth}
はば\qty{8}{m}の屋根を、左が\ang{30}、右が\ang{45}の、左右ちがう傾きで作ります。上の角は $\ang{180}-\ang{30}-\ang{45}=\ang{105}$ です。正弦定理で、ななめの部材の長さは、
\bkivhline{\begin{tabular}{@{}l@{}}左 $=\dfrac{8\sin\ang{45}}{\sin\ang{105}}\approx\qty{5.86}{m}$\\[4pt] 右 $=\dfrac{8\sin\ang{30}}{\sin\ang{105}}\approx\qty{4.14}{m}$\end{tabular}}
実際の家の屋根は、左右同じ5寸こう配（傾き0.5、約\ang{26.6}）などが多いのですが、どんな形でも計算のしかたは同じです。三角形は形が変わらない（4-1）ので、屋根や橋の骨組みは三角形を組み合わせて作ります。部材を切る長さは、このような計算で決めます。
\end{minipage}\hfill
\begin{minipage}[c]{0.36\linewidth}\centering
\begin{tikzpicture}[x=1mm, y=1mm]
  \coordinate (L) at (0,0); \coordinate (Rr) at (44,0); \coordinate (T) at ({32.2*cos(30)},{32.2*sin(30)});
  \fill[bkOrange!12] (L) -- (Rr) -- (T) -- cycle;
  \draw[bkOrange!70!black, line width=1.6pt] (L) -- (Rr) -- (T) -- cycle;
  \draw[bkOrange!70!black, line width=0.9pt] (T) -- ({32.2*cos(30)},0);
  \draw[bkInk, line width=0.5pt] (6,0) arc (0:30:6);
  \node[font=\sffamily\footnotesize] at (15:9) {\ang{30}};
  \draw[bkInk, line width=0.5pt] ($(Rr)+(-5,0)$) arc (180:135:5);
  \node[font=\sffamily\footnotesize] at ($(Rr)+(157:8.2)$) {\ang{45}};
  \node[font=\sffamily\footnotesize, anchor=north] at (22,-1.2) {\qty{8}{m}};
\end{tikzpicture}
\end{minipage}
\end{butsuri}

\begin{butsuri}
\par\addvspace{3pt}\textbf{三角測量：日本の地図を作る}\par\noindent
\begin{minipage}[c]{0.6\linewidth}
明治時代から、日本の地図は三角測量で作られてきました。山の上などに「三角点」を置き、2つの三角点のあいだの長さ（基線）を、とても正確にはかります。あとは、三角形の角だけをはかれば、正弦定理で次々に辺の長さが求められます。角は遠くからでもはかれるので、山や川をこえて、三角形を国じゅうにつないでいけるのです。今は人工衛星を使った測量も使われていますが、三角点は今も全国にあります。
\end{minipage}\hfill
\begin{minipage}[c]{0.36\linewidth}\centering
\begin{tikzpicture}[x=1mm, y=1mm]
  \coordinate (P1) at (0,0); \coordinate (P2) at (16,2); \coordinate (P3) at (8,15); \coordinate (P4) at (26,16); \coordinate (P5) at (18,29); \coordinate (P6) at (36,28);
  \fill[bkGreen!12] (P1) -- (P2) -- (P4) -- (P6) -- (P5) -- (P3) -- cycle;
  \draw[bkInk!70, line width=0.6pt] (P1) -- (P2) -- (P3) -- cycle (P2) -- (P4) -- (P3) (P3) -- (P5) -- (P4) (P4) -- (P6) -- (P5);
  \draw[bkRed, line width=1.6pt] (P1) -- (P2);
  \foreach \p in {P1,P2,P3,P4,P5,P6} {\fill[bkInk] (\p) circle (0.8);}
  \node[font=\sffamily\footnotesize, text=bkRed!80!black, anchor=north] at (8,0.3) {基線};
\end{tikzpicture}
\end{minipage}

\par\addvspace{4pt}\textbf{クレーンの「ブーム」を支えるワイヤ}\par\noindent
\begin{minipage}[c]{0.6\linewidth}
クレーンの長い腕（ブーム）は、根もとのピンを中心に回り、先をワイヤで引いて角度を変えます。ブームが\qty{30}{m}、ピンの真上\qty{10}{m}の所からワイヤを引くとき、ブームを水平から\ang{60}に起こすと、ブームとピンの真上へのばした線の角は\ang{30}です。余弦定理で（$30^2=900$、$10^2=100$、$2\times30\times10=600$）、ワイヤの長さは
\bkivhline{$\sqrt{900+100-600\cos\ang{30}}\approx\qty{21.9}{m}$}
ワイヤを巻き取って短くすると、ブームが起き上がります。
\end{minipage}\hfill
\begin{minipage}[c]{0.36\linewidth}\centering
\begin{tikzpicture}[x=1mm, y=1mm]
  \fill[bkGray!40] (-8,-3) rectangle (14,0);
  \fill[bkYellow!80!black] (-6,0) rectangle (6,6);
  \coordinate (Pin) at (0,6); \coordinate (Top) at (0,16);
  \coordinate (Tip) at ($(Pin)+(60:30)$);
  \draw[bkInk, line width=0.8pt] (Pin) -- (Top);
  \draw[bkYellow!70!black, line width=2.6pt, line cap=round] (Pin) -- (Tip);
  \draw[bkInk, line width=0.7pt] (Top) -- (Tip);
  \draw[bkInk, line width=0.5pt] ($(Pin)+(0,5)$) arc (90:60:5);
  \draw[bkInk, line width=0.4pt] ($(Pin)+(75:5)$) -- ($(Pin)+(-5,3.4)$);
  \node[font=\sffamily\footnotesize, anchor=south east, inner sep=1pt] at ($(Pin)+(-5,1.8)$) {\ang{30}};
  \draw[bkInk, line width=0.6pt] (Tip) -- ++(0,-8);
  \fill[bkInk] ($(Tip)+(-1.2,-10)$) rectangle ($(Tip)+(1.2,-8)$);
  \node[font=\sffamily\footnotesize, anchor=east] at ($(Pin)+(-0.6,8.6)$) {10};
  \node[font=\sffamily\footnotesize, anchor=north west] at ($(Pin)!0.55!(Tip)+(0.8,-0.5)$) {30};
  \node[font=\sffamily\footnotesize, anchor=south east] at ($(Top)!0.5!(Tip)+(-0.3,0.3)$) {ワイヤ};
\end{tikzpicture}
\end{minipage}
\end{butsuri}

\begin{column}{スマートフォンの位置も、三角形で決める}
スマートフォンの地図アプリは、人工衛星や Wi-Fi の電波から、「Aから何m」「Bから何m」という距離を求めて、今いる場所を決めています。2つの基地局A、Bのあいだの距離と、それぞれからの距離がわかると、3辺のわかった三角形ができます。たとえば $\mathrm{AB}=\qty{100}{m}$、Aから\qty{80}{m}、Bから\qty{60}{m}なら、余弦定理で
\bkivhline{$\cos A=\dfrac{100^2+80^2-60^2}{2\times100\times80}=\dfrac{\num{12800}}{\num{16000}}=0.8$}
Aの角は約\ang{36.9}とわかり、AからBを見た向きから\ang{36.9}回った向きに\qty{80}{m}の所にいる、と決まります（ABのどちらがわかは、3つめの基地局で決める）。三角形を解く計算は、毎日あなたのポケットの中でも行われているのです。
\end{column}

\begin{column}{星までの距離も、三角形ではかる}
地球は太陽のまわりを回っているので、半年たつと、太陽をはさんで反対がわ（約3億km先）に移ります。そのあいだに、近くの星は、遠くの星を背景にして、ほんの少しだけ位置がずれて見えます。地球と太陽の距離（約1.5億km）を底辺にした、とても細長い三角形ができるのです。太陽系にいちばん近い星プロキシマ・ケンタウリでは、このずれの角（年周視差）は約0.77秒（1秒は\ang{1}の3600分の1）しかありません。それでも $\tan$ を使えば、距離は約40兆km、光で約4.2年かかる遠さだと計算できます。三角形の計算は、宇宙のものさしにもなっています。
\end{column}

%% ---------------- 章末問題 ----------------
\bkendproblems
\bkproblemgroup{基本}{3}

\bkproblem 三角形ABCの面積を求めましょう。\quad（1）$a=7$、$b=4$、$C=\ang{30}$\quad（2）$b=6$、$c=5$、$A=\ang{120}$\mbox{（\bkreidairef{1}）}
\bkpans{（1）\nobreak $\dfrac{1}{2}\times7\times4\times\dfrac{1}{2}=7$。（2）\nobreak $\dfrac{1}{2}\times6\times5\times\dfrac{\sqrt3}{2}=\dfrac{15\sqrt3}{2}$（約13.0）。}

\bkproblem 三角形ABCで、$a=6$、$A=\ang{30}$、$B=\ang{45}$ です。$b$ と、外接円の半径 $R$ を求めましょう。\mbox{（\bkreidairef{2}）}
\bkpans{$b=\dfrac{6\sin\ang{45}}{\sin\ang{30}}=6\sqrt2$（約8.49）。$2R=\dfrac{6}{\sin\ang{30}}=12$ より $R=6$。}

\bkproblem $\mathrm{AB}=\qty{60}{m}$、Aから見た角が\ang{50}、Bから見た角が\ang{60}の三角形ABCで、ACとBCの長さを小数第1位まで求めましょう。\mbox{（\bkreidairef{3}）}
\bkpans{$C=\ang{70}$。$\mathrm{AC}=\dfrac{60\sin\ang{60}}{\sin\ang{70}}\approx\qty{55.3}{m}$。$\mathrm{BC}=\dfrac{60\sin\ang{50}}{\sin\ang{70}}\approx\qty{48.9}{m}$。}

\bkproblem 三角形ABCで、$c$ を求めましょう。\quad（1）$a=8$、$b=5$、$C=\ang{60}$\quad（2）$a=2$、$b=3$、$C=\ang{120}$\mbox{（\bkreidairef{4}）}
\bkpans{（1）\nobreak $c^2=64+25-40=49$、$c=7$。（2）\nobreak $c^2=4+9+6=19$、$c=\sqrt{19}$（約4.36）。}

\bkproblem 3辺が $a=7$、$b=5$、$c=8$ の三角形ABCで、角 $A$ を求めましょう。\mbox{（\bkreidairef{5}）}
\bkpans{$\cos A=\dfrac{25+64-49}{2\times5\times8}=\dfrac{40}{80}=\dfrac{1}{2}$。$A=\ang{60}$。}

\bkproblem 池をはさんだ2点A、Bについて、点Cで $\mathrm{CA}=\qty{50}{m}$、$\mathrm{CB}=\qty{70}{m}$、Cの角\ang{40}をはかりました。ABを小数第1位まで求めましょう。\mbox{（\bkreidairef{6}）}
\bkpans{$c^2=7400-7000\cos\ang{40}\approx2037.7$、$c\approx\qty{45.1}{m}$。}

\bkproblem \qty{20}{N}と\qty{30}{N}の力が、あいだの角\ang{60}で1つの物を引いています。合力の大きさを小数第1位まで求めましょう。\mbox{（\bkreidairef{9}）}
\bkpans{$R^2=400+900-1200\cos\ang{120}=1900$、$R\approx\qty{43.6}{N}$。}

\bkproblemgroup{チャレンジ}{3}

\bkproblem $L_1=50$、$L_2=40$（cm）の2リンクのうでで、$\theta_1=\ang{45}$、$\theta_2=\ang{90}$ のとき、手の座標を小数第1位まで求めましょう。\mbox{（\bkreidairef{7}）}
\bkpans{$\theta_1+\theta_2=\ang{135}$ より、$x\approx35.36-28.28\approx7.1$、$y\approx35.36+28.28\approx63.6$。}

\bkproblem $L_1=40$、$L_2=30$（cm）の2リンクのうでで、手を $(30,\ 50)$ に置きます。ひじを下にする形の $\theta_2$ と $\theta_1$ を小数第1位まで求めましょう。\mbox{（\bkreidairef{8}）}
\bkpans{$d\approx58.31$。$\cos\angle\mathrm{OEP}=-0.375$ より $\theta_2\approx\ang{68.0}$。$\alpha\approx\ang{59.04}$、$\cos\beta\approx0.8789$ より $\beta\approx\ang{28.49}$。$\theta_1\approx\ang{30.5}$。}

\bkproblem 重力 $W=\qty{49}{N}$ の物を2本のロープでつるします。ロープが真上からかたむいた角が $\theta_{\mathrm{L}}=\theta_{\mathrm{R}}=\ang{45}$ のとき、$T_{\mathrm{L}}$、$T_{\mathrm{R}}$ を小数第1位まで求めましょう。\mbox{（\bkreidairef{10}）}
\bkpans{$\theta_{\mathrm{L}}+\theta_{\mathrm{R}}=\ang{90}$。$T_{\mathrm{L}}=T_{\mathrm{R}}=\dfrac{49\times\sin\ang{45}}{\sin\ang{90}}\approx\qty{34.6}{N}$。}

\begin{furikaeri}
\begin{checklist}
\item 鈍角の三角比を使い、三角形の面積を $\dfrac{1}{2}ab\sin C$ で求められる。\bkfurref{bk:4-7:menseki}
\item 正弦定理と余弦定理を使い分けて、三角形の辺や角を求められる。\bkfurref{bk:4-7:seigen}
\item ロボットのうでの関節の角を、余弦定理で求められる。\bkfurref{bk:4-7:robot}
\end{checklist}
\end{furikaeri}

\bkfukidashi{mina}{行けない場所までの距離も、ロボットのうでの角も、三角形で計算できたね。次の章では、回る点の動きを三角関数で表すよ。}
